% Source presentations for register row N-c-010.

\subsection{Three tractability mechanisms and the normative OWL boundary}
\label{sec:p09-presentation-n-c-010}

\subsubsection{The \(\mathcal{EL}^{++}\) saturation of Baader--Brandt--Lutz}

\paragraph{Sources and proof boundary.}
The principal source is the six-page official IJCAI paper, printed
pp.~364--369.  It expressly sends proofs and technical details to the authors'
thirty-two-page LTCS-05-01 report.
IJCAI numbering controls statements below; report numbering identifies
report-only material.  The procedure here uses the source's finite completion
state and rules with three explicit handbook corrections.  A conservative
concrete-atom abbreviation prepass gives linear written normalization size
also for unbounded concrete arities.  A fresh nominal then marks the queried
concept before saturation, and concrete-atom propagation checks that all
required features are already defined in the relevant concrete domain.
The soundness and countermodel arguments below prove this corrected
procedure for the full displayed language, including nominals and partial
features.  The source's unmarked simultaneous-classification claim is not used
as a premise of that proof.  The official proceedings reserve reproduction
rights, so the handbook reconstructs the mathematics rather than reproducing
the printed pages
\cite{BaaderBrandtLutz2005,BaaderBrandtLutz2005TR}.

\paragraph{The language, concrete domains, and CBoxes.}
Let \(N_C,N_R,N_I,N_F\) be the source sets of concept names, role names,
individual names, and feature names.  A concrete domain is
\(\mathcal D=(\Delta^{\mathcal D},P^{\mathcal D})\), where every
\(p\in P^{\mathcal D}\) has a positive arity and an extension over
\(\Delta^{\mathcal D}\).  When
\(\mathcal D_1,\ldots,\mathcal D_n\) are used together, their carriers are
pairwise disjoint.  The source concepts are exactly
\[
 C,D::=A\mid\top\mid\bot\mid\{a\}\mid C\sqcap D
       \mid\exists r.C\mid p(f_1,\ldots,f_k),
\]
where \(A\in N_C\), \(a\in N_I\), \(r\in N_R\), the \(f_i\in N_F\), and
\(p\in P^{\mathcal D_j}\) has arity \(k\).  The six non-name forms stand for
top, bottom, a nominal, conjunction, an existential role restriction, and a
concrete-domain atom; none is OWL syntax imported into the paper.

An interpretation \(\mathcal I=(\Delta^{\mathcal I},\cdot^{\mathcal I})\)
has \(\Delta^{\mathcal I}\ne\varnothing\) and assigns
\[
\begin{gathered}
 A^{\mathcal I}\subseteq\Delta^{\mathcal I},\qquad
 r^{\mathcal I}\subseteq\Delta^{\mathcal I}\times\Delta^{\mathcal I},
 \qquad a^{\mathcal I}\in\Delta^{\mathcal I},\\
 f^{\mathcal I}:\Delta^{\mathcal I}\rightharpoonup
       \bigcup_{j=1}^{n}\Delta^{\mathcal D_j}.
\end{gathered}
\]
Besides the usual clauses for \(\top,\bot,\sqcap\),
\[
\begin{aligned}
 \{a\}^{\mathcal I}&=\{a^{\mathcal I}\},\\
 (\exists r.C)^{\mathcal I}
 &=\{x:\exists y\,((x,y)\in r^{\mathcal I}\land y\in C^{\mathcal I})\},\\
 p(f_1,\ldots,f_k)^{\mathcal I}
 &=\{x:\exists y_1,\ldots,y_k\in\Delta^{\mathcal D_j}
       (\bigwedge_{i=1}^{k}f_i^{\mathcal I}(x)=y_i
        \land (y_1,\ldots,y_k)\in p^{\mathcal D_j})\}.
\end{aligned}
\]
A concrete atom is consequently false when a required feature is undefined.

A \emph{constraint box}, or CBox, is a finite set whose members may be GCIs
\(C\sqsubseteq D\) and RIs
\(r_1\circ\cdots\circ r_k\sqsubseteq r\).  The former require
\(C^{\mathcal I}\subseteq D^{\mathcal I}\); the latter require relational
composition on the left to be included in the single role on the right.  A
finite GCI-only set is a general TBox in this source.  Hence a CBox and a
general TBox are not synonyms.  For a CBox \(\mathcal C\),
\[
 C\sqsubseteq_{\mathcal C}D
 \quad\Longleftrightarrow\quad
 C^{\mathcal I}\subseteq D^{\mathcal I}
 \text{ for every model \(\mathcal I\) of \(\mathcal C\)}.
\]
There is no global unique-name assumption: selected inequality can instead be
forced by \(\{a\}\sqcap\{b\}\sqsubseteq\bot\)
\cite[table~1 and \S2]{BaaderBrandtLutz2005}.

\paragraph{Normal form and the written-size convention.}
\label{sec:p09-elpp-normalization}
For any finite CBox \(\mathcal C\), its basic descriptions form exactly
the smallest set \(\mathsf{BC}_{\mathcal C}\) containing \(\top\), all
concept names occurring in \(\mathcal C\), and all nominals and concrete
atoms occurring as concept subdescriptions in \(\mathcal C\).
There are no other elements: \(\bot\), conjunctions and existential
restrictions are not basic.  The same definition is used for each later
normalized or marked CBox.

In this normalization discussion, the source-local notation
\(|\mathcal C|\) means the number of symbols in a fully written CBox,
not its number of axioms
\cite[\S3.1, before lem.~2]{BaaderBrandtLutz2005TR}.
Names are atomic vocabulary symbols in this measure.  Every occurrence
of a connective, role, feature argument and punctuation symbol is counted;
a concrete atom \(p(f_1,\ldots,f_k)\) has length proportional to \(k+1\),
rather than unit length.  Write \(\ell(E)\) for the corresponding written
length of a concept description \(E\).
No upper bound on concrete-predicate arities is assumed.
The later bars in \(b=|\mathsf{BC}_{\mathcal K}|\) and
\(m_R=|R_{\mathcal K}|\) mean ordinary cardinalities of the indicated
finite vocabularies.  Encoded complexity, which also counts written
identifiers and concrete-predicate parameters, is addressed after the
normalization proof.

A normalized CBox has only
\[
\begin{array}{c@{\qquad}c}
 C_1\sqsubseteq D & C_1\sqcap C_2\sqsubseteq D\\
 C_1\sqsubseteq\exists r.C_2 & \exists r.C_1\sqsubseteq D
\end{array}
\qquad
\begin{array}{c}
 r\sqsubseteq s\\
 r_1\circ r_2\sqsubseteq s,
\end{array}
\]
where \(C_1,C_2\in\mathsf{BC}_{\mathcal C}\) and
\(D\in\mathsf{BC}_{\mathcal C}\cup\{\bot\}\).

\paragraph{A conservative concrete-atom abbreviation step.}
The following prepass is a handbook addition to the report's NF schedule.
For each syntactically distinct occurring concrete atom \(X\), introduce a
fresh concept name \(A_X\), with different atoms receiving different names.
Its subscript identifies the associated atom
in this description of the construction; the output uses an ordinary fresh
name, not a copy of the expression \(X\) inside that name.
Replace every occurrence of \(X\) in the input GCIs by \(A_X\), and
adjoin the two defining GCIs
\[
 A_X\sqsubseteq X,\qquad X\sqsubseteq A_X.
 \tag{N0}
\]
Role inclusions are unchanged.  Denote the replaced input by
\(\mathcal C^\flat\) and the set of these defining GCIs by
\(\mathcal E\).  For an original description \(E\), write \(E^\flat\)
for the result of the same atom replacement within \(E\).
The definitions in \(\mathcal E\) are already normalized and are
protected: the prepass is not applied again to their atoms, and subsequent
NF rewrites act only on \(\mathcal C^\flat\).
Their full atom expressions remain in the final output.
The basic set of \(\mathcal C^\flat\cup\mathcal E\) includes both
\(A_X\) and \(X\), by the exact definition above.

Every original interpretation \(\mathcal I\) has an expansion
\(\mathcal I^+\) satisfying \(\mathcal E\), by setting
\(A_X^{\mathcal I^+}=X^{\mathcal I}\) for every such atom and keeping all
original symbols unchanged.  This assignment uses the given partial-feature
semantics of \(X\), including undefined-feature cases.
A structural induction on an original concept description \(E\) then gives
\(E^{\mathcal I}=(E^\flat)^{\mathcal I^+}\): the concrete-atom case is
the defining assignment, the other atomic cases are unchanged, conjunction
uses intersection, and existential restriction uses the unchanged role and
the inductive equality of its filler.
Thus an original model expands to a model of
\(\mathcal C^\flat\cup\mathcal E\).
Conversely, the two inclusions in (N0) force
\(A_X^{\mathcal J}=X^{\mathcal J}\) in any model \(\mathcal J\)
of the expanded CBox.  The same induction equates original and replaced
descriptions, so its reduct satisfies every original GCI; the RIs were
unchanged.  The prepass is therefore a model-conservative extension.

Put \(N=|\mathcal C|\).
Concrete atoms contain feature names rather than nested concepts, so their
occurrence subexpressions are disjoint in the written input.
Consequently the sum of \(\ell(X)\) over distinct occurring atoms is at
most the sum over all occurrences, and hence at most \(N\).
The two definitions for \(X\) have total length
\(2\ell(X)+O(1)\); their total length is \(O(N)\).
The replacement part has length \(O(N)\) as well, because every occurrence
is replaced by a single name.  This proves linear written size of the
prepass output.  It neither treats a whole atom as a unit-cost expression
nor assumes a shared-expression final output.
Distinct atoms can be identified by finite syntactic comparison; even the
direct pairwise comparison procedure takes only polynomial time.

\paragraph{The seven source normalization transformations.}
The conference paper states the reduction; the report supplies the following
NF1--NF7 rules
\cite[lem.~1]{BaaderBrandtLutz2005}
\cite[fig.~1 and \S3.1]{BaaderBrandtLutz2005TR}.
In these schemas, \(B,C,D,E\) range over arbitrary concept descriptions,
not just concept names.  A hat requires that the indicated description
is not basic in the current CBox; in particular \(\bot\) is nonbasic.
The letters \(r_i,r,s\) denote role names.
Every \(A\) introduced by a rule is a fresh concept name, and every \(u\)
is a fresh role name, distinct also from all prepass names and earlier
introduced names.  Fresh names enter the basic vocabulary when they are
introduced.
\begin{enumerate}
\item[(NF1)] for \(k>2\), replace
\(r_1\circ\cdots\circ r_k\sqsubseteq s\) by
\(r_1\circ\cdots\circ r_{k-1}\sqsubseteq u\) and
\(u\circ r_k\sqsubseteq s\);
\item[(NF2)] replace \(C\sqcap\widehat D\sqsubseteq E\) by
\(\widehat D\sqsubseteq A\) and \(C\sqcap A\sqsubseteq E\), modulo
commutativity of conjunction;
\item[(NF3)] replace \(\exists r.\widehat C\sqsubseteq D\) by
\(\widehat C\sqsubseteq A\) and \(\exists r.A\sqsubseteq D\);
\item[(NF4)] delete \(\bot\sqsubseteq D\);
\item[(NF5)] replace \(\widehat C\sqsubseteq\widehat D\) by
\(\widehat C\sqsubseteq A\) and \(A\sqsubseteq\widehat D\);
\item[(NF6)] replace \(B\sqsubseteq\exists r.\widehat C\) by
\(B\sqsubseteq\exists r.A\) and \(A\sqsubseteq\widehat C\);
\item[(NF7)] replace \(B\sqsubseteq C\sqcap D\) by
\(B\sqsubseteq C\) and \(B\sqsubseteq D\).
\end{enumerate}
Exhaust NF1--NF4 on the mutable part first, then exhaust NF5--NF7.
Keep \(\mathcal E\) in the CBox throughout, without rewriting its
already-normal definitions.  This is the source's two-phase order,
applied after the additional prepass.

Each NF step is conservative.  For NF1, expand an original model by
interpreting \(u\) as \(r_1\circ\cdots\circ r_{k-1}\).
Conversely, inclusion of that prefix in \(u\), followed by
\(u\circ r_k\sqsubseteq s\), implies the old RI.
For NF2 and NF3, interpret the new \(A\) as the extension of the
extracted \(\widehat D\) and \(\widehat C\), respectively.
In a model of the replacement axioms, those extracted descriptions are
included in \(A\); monotonicity of intersection and existential restriction
then gives the original GCI.
For NF5, interpret \(A\) as \(\widehat C\); conversely transitivity
of inclusion recovers \(\widehat C\sqsubseteq\widehat D\).
For NF6, interpret \(A\) as \(\widehat C\); conversely
\(A\sqsubseteq\widehat C\) gives
\(\exists r.A\sqsubseteq\exists r.\widehat C\).
NF4 deletes a GCI with empty antecedent extension, and NF7 uses exactly
the equivalence between inclusion in an intersection and the two separate
inclusions.  Freshness leaves every other axiom, including \(\mathcal E\),
unchanged in each expansion.  Finite composition of these expansion and
reduct arguments proves conservativity of the whole normalization.

\paragraph{Termination, application count and linear written size.}
The count here is proved for the amended construction, rather than inferred
from the number of axioms in the input.
For the mutable prepass output, let \(g\) be its number of GCIs,
\(e\) the number of concept-child edges in all its left-hand concept
trees, and
\[
 h=\sum_{\text{input RIs }r_1\circ\cdots\circ r_k\sqsubseteq s}
            \max(k-2,0).
\]
Let \(a\) and \(v\) be, respectively, the numbers of conjunction and
existential-constructor occurrences in its right-hand concept trees.
All five counts are \(O(N)\).

An NF2 or NF3 application cuts one existing left-tree edge whose child
is nonbasic, replacing that child occurrence by a fresh basic name and
moving the original subtree to the left of a new GCI.
No old concept subtree is copied in phase one.  The replacement edge has
a basic child and cannot be cut again.
Charging each step to its removed original edge therefore bounds the
combined number of NF2 and NF3 applications by \(e\).
At most \(g+e\) left-hand GCI roots are ever created in this phase, so
NF4 applies at most \(g+e\) times.  Each NF1 step shortens its one
remaining long chain by one and emits a chain of length two; hence it
applies exactly \(h\) times.
Phase one therefore takes at most
\[
 g+2e+h
\]
steps and leaves at most \(g+e\) mutable GCIs.

After phase one, every mutable left side is basic, a conjunction of two
basic descriptions, or an existential restriction with basic filler;
a left side \(\bot\) has been deleted.
There are no concrete atoms in the mutable part, since the prepass
replaced them all.  Its basic leaves are therefore names, nominals and
\(\top\), all of bounded symbol length.  Every mutable left side has
bounded symbol length independent of the input's predicate arities.

Phase two never copies a right-hand concept subtree.
NF7 removes its conjunction root and places the two children on separate
right sides, so it can be charged to one of the \(a\) original right-hand
conjunction occurrences.
NF6 extracts its nonbasic filler while replacing that occurrence by a
fresh name.  The same existential occurrence then has a basic filler and
cannot trigger NF6 again, so there are at most \(v\) such steps.
NF5 moves its right side without copying it and replaces the corresponding
nonbasic left side, for further right-side processing, by a fresh name.
To count these steps, count GCIs with both sides nonbasic.  Initially in
phase two there are at most \(g+e\) such GCIs.
NF5 reduces this count by one; NF6 does not increase it; NF7 increases
it by at most one.
Thus NF5 applies at most \(g+e+a\) times.
The entire NF schedule uses at most
\[
 2g+3e+h+v+2a=O(N)
 \tag{N1}
\]
applications.  In particular it terminates.

Phase two keeps each left side equal to a phase-one left side or a fresh
name, with possible copies under NF7.  Those left sides all have bounded
length.  NF1--NF3 and NF5--NF6 increase written length by only a fixed
constant per step: existing subexpressions are moved, and only a bounded
number of fresh-name and punctuation occurrences is added.
NF4 only deletes text.  NF7 copies only its bounded-length left side and
adds bounded punctuation.  Combining this observation with (N1) and the
prepass bound proves that the fully written final CBox has size \(O(N)\),
including both full copies of every distinct concrete atom in \(\mathcal E\).

At termination no phase-one rule has become newly applicable: phase two
introduces only fresh basic left sides or copies of the already shallow
ones.  A nonbasic left side and nonbasic right side would still trigger
NF5.  With a basic left side, a conjunction on the right triggers NF7,
and an existential with nonbasic filler triggers NF6.
All remaining mutable GCIs therefore have one of the four required
normal forms; the protected definitions and the RIs do as well.

For example, write \(X_k=p_k(f_1,\ldots,f_k)\).
The input \(X_k\sqsubseteq A_1\sqcap\cdots\sqcap A_k\) becomes
\(A_{X_k}\sqsubseteq A_1\sqcap\cdots\sqcap A_k\) together with the
two definitions of \(A_{X_k}\).
NF7 produces \(A_{X_k}\sqsubseteq A_i\) for each \(i\).
The full \(k\)-argument atom occurs only twice, in its protected definitions,
so the complete output has linear symbol length.  The same construction
handles repeated occurrences by using one definition pair per distinct atom.

For effective complexity, let \(L_{\mathrm{enc}}\) count a direct encoding
of the fully written input, including identifiers and concrete-predicate parameters.
Fresh names may be encoded by consecutive fresh indices of polynomial bit
length.  There are only linearly many output symbol occurrences in the
source measure, and each original name or parameter has length at most
\(L_{\mathrm{enc}}\).  Hence even a direct written encoding of the output
has polynomial length in \(L_{\mathrm{enc}}\).
Finite comparison for the prepass and direct rescanning for NF applications
also take polynomial time.  The linear claim just proved is a written
symbol-size and rewrite-count claim; the effective polynomial bound does
not assume unit-cost arithmetic on unbounded predicate parameters.

\paragraph{Arbitrary queries before normalization.}
For a raw query \(C\sqsubseteq_{\mathcal C}D\), first adjoin
\(A_0\sqsubseteq C\) and \(D\sqsubseteq B_0\), with two fresh concept
names.  Apply the concrete-atom prepass to this entire expanded input,
including atoms from the query, and then apply the proved two-phase
normalization.
If the original subsumption holds, every element of \(A_0\) in a
model of the expanded raw CBox belongs successively to \(C,D,B_0\).
Conversely, an original countermodel with
\(x\in C^{\mathcal I}\setminus D^{\mathcal I}\) expands by
\(A_0^{\mathcal I}=C^{\mathcal I}\) and
\(B_0^{\mathcal I}=D^{\mathcal I}\), retaining that counterexample.
The prepass and NF transformations have the expansion/reduct properties
proved above and do not change these queried names' extensions.
Thus the atomic query is preserved in both directions, and the normalized
output has size linear in
\(|\mathcal C|+\ell(C)+\ell(D)+1\) under the stated source convention.

\paragraph{Marking the queried element.}
\label{sec:p09-elpp-query-marking}
Let \(\mathcal C\) now be a normalized CBox and let
\(A_0\sqsubseteq B_0\) be the atomic query.  Take an individual name
\(a_\ast\) fresh for \(\mathcal C,A_0,B_0\), and put
\[
 N_\ast=\{a_\ast\},\qquad
 \mathcal K=\mathcal C\cup
       \{N_\ast\sqsubseteq A_0,\ B_0\sqsubseteq B_0\}.
\]
The tautology ensures that the queried name \(B_0\) occurs in the finite
vocabulary even when it did not occur in \(\mathcal C\).
The new axioms are already normalized.  Then
\[
 A_0\sqsubseteq_{\mathcal C}B_0
 \quad\Longleftrightarrow\quad
 N_\ast\sqsubseteq_{\mathcal K}B_0.
 \tag{Q}
\]
For the forward implication, a \(\mathcal K\)-model is a
\(\mathcal C\)-model and its element \(a_\ast^{\mathcal I}\) belongs to
\(A_0^{\mathcal I}\).  For the reverse implication, any
\(\mathcal C\)-model with
\(x\in A_0^{\mathcal I}\setminus B_0^{\mathcal I}\) expands to a
\(\mathcal K\)-model by setting \(a_\ast^{\mathcal I}=x\).
Freshness is essential to this expansion.  In particular, the marking is a
query reduction, not a model-conservative extension: if \(A_0\) is forced
empty, \(\mathcal K\) has no model and the original subsumption is true.
The saturation below is run on \(\mathcal K\), and its queried label is
\(S(N_\ast)\).

\paragraph{The finite saturation state and its abbreviations.}
For the marked normalized CBox \(\mathcal K\), let \(R_{\mathcal K}\) be its finite role
vocabulary.  The algorithm computes
\[
\begin{aligned}
 S&:\mathsf{BC}_{\mathcal K}\longrightarrow
   \mathcal P(\mathsf{BC}_{\mathcal K}\cup\{\top,\bot\}),\\
 R&:R_{\mathcal K}\longrightarrow
   \mathcal P(\mathsf{BC}_{\mathcal K}\times
              \mathsf{BC}_{\mathcal K}),
\end{aligned}
\]
initialized by \(S(C)=\{C,\top\}\) and \(R(r)=\varnothing\).  The intended
invariants are
\[
 D\in S(C)\Rightarrow C\sqsubseteq_{\mathcal K}D,
 \qquad
 (C,D)\in R(r)\Rightarrow C\sqsubseteq_{\mathcal K}\exists r.D.
\]
The abbreviation \(C\rightsquigarrow D\) holds exactly when a finite nonempty
sequence \(C_1,\ldots,C_k\), with \(k\geq1\), ends in \(D\), begins either in \(C\) or in some nominal
\(\{b\}\), and each adjacent pair belongs to \(R(r_j)\) for some role
\(r_j\).  It is computed from the current \(R\); it is not a new interpreted
role.

For a set \(\Gamma\) of concepts, put
\[
 \operatorname{con}_j(\Gamma)=
 \bigwedge_{\substack{p(f_1,\ldots,f_k)\in\Gamma\\
                       p\in P^{\mathcal D_j}}}
 p(f_1,\ldots,f_k).
\]
Let \(F_j(C)\) be the finite set of feature names occurring in
\(\operatorname{con}_j(S(C))\).  A solution of this conjunction is a map
from exactly these feature names to \(\Delta^{\mathcal D_j}\) satisfying
each displayed atom.  The empty conjunction has the empty map as a solution.
Implication to an atom all of whose features lie in \(F_j(C)\) means truth
under every such solution.  These are finite concrete constraints; they do
not assign values to features not required at the abstract element.

A concrete domain is \emph{\(p\)-admissible} when satisfiability and
implication for finite conjunctions of its atoms are decidable in polynomial
time, and it is convex: whenever such a conjunction implies a nonempty finite
disjunction of atoms, it implies a disjunct.  The finite supports and
empty-conjunction convention fix the tests used here.  Polynomial constraint
reasoning makes CR7--CR8 executable; convexity supplies the simultaneous
avoiding assignment in the completeness proof.

\paragraph{Partial-feature propagation.}
\label{sec:p09-elpp-partial-features}
CR8 below requires every feature of its conclusion to belong to \(F_j(C)\).
The positive atoms already in \(S(C)\) then ensure that, at any instance of
\(C\), each required feature is defined and takes a value in
\(\Delta^{\mathcal D_j}\).  Without that condition, a concrete implication
under total variable assignments need not be an implication between concepts
with partial features.  For example, the source rational domain has a
universal unary predicate \(\top_{\mathbb Q}\).  In a one-element abstract
model with \(A\) true, \(B\) false and \(f\) undefined, the normalized GCI
\(\top_{\mathbb Q}(f)\sqsubseteq B\) holds.  Consequently it does not force
\(A\sqsubseteq B\).  The guarded rule leaves
\(\top_{\mathbb Q}(f)\) absent from that otherwise unconstrained \(A\)-label:
the feature \(f\) is not in \(F_{\mathbb Q}(A)\).
This feature-support condition is part of the corrected procedure; the source
Table~2 states its concrete implication test without that condition
\cite[table~1, table~2 and prop.~5]{BaaderBrandtLutz2005}.

\paragraph{All eleven completion rules.}
Each applicable rule makes the one displayed monotone addition; none branches
or removes information.
\begin{enumerate}
\item[(CR1)] If \(C'\in S(C)\), \(C'\sqsubseteq D\in\mathcal K\), and
\(D\notin S(C)\), add \(D\) to \(S(C)\).
\item[(CR2)] If \(C_1,C_2\in S(C)\),
\(C_1\sqcap C_2\sqsubseteq D\in\mathcal K\), and \(D\notin S(C)\), add
\(D\) to \(S(C)\).
\item[(CR3)] If \(C'\in S(C)\),
\(C'\sqsubseteq\exists r.D\in\mathcal K\), and \((C,D)\notin R(r)\), add
\((C,D)\) to \(R(r)\).
\item[(CR4)] If \((C,D)\in R(r)\), \(D'\in S(D)\),
\(\exists r.D'\sqsubseteq E\in\mathcal K\), and \(E\notin S(C)\), add
\(E\) to \(S(C)\).
\item[(CR5)] If \((C,D)\in R(r)\), \(\bot\in S(D)\), and
\(\bot\notin S(C)\), add \(\bot\) to \(S(C)\).
\item[(CR6)] If \(\{a\}\in S(C)\cap S(D)\),
\(C\rightsquigarrow D\), and \(S(D)\nsubseteq S(C)\), add all of \(S(D)\)
to \(S(C)\).
\item[(CR7)] If \(\operatorname{con}_j(S(C))\) is unsatisfiable in
\(\mathcal D_j\) and \(\bot\notin S(C)\), add \(\bot\) to \(S(C)\).
\item[(CR8)] If \(X=p(f_1,\ldots,f_k)\in\mathsf{BC}_{\mathcal K}\),
\(p\in P^{\mathcal D_j}\), every \(f_i\) belongs to \(F_j(C)\),
\(\operatorname{con}_j(S(C))\) implies \(X\) in \(\mathcal D_j\), and
\(X\notin S(C)\), add \(X\) to \(S(C)\).
\item[(CR9)] If atoms from distinct concrete domains occur in \(S(C)\) and
share a feature name, add \(\bot\) to \(S(C)\) if it is absent.  Formally the
premises are \(p(f_1,\ldots,f_k),p'(f'_1,\ldots,f'_{k'})\in S(C)\),
\(p\in P^{\mathcal D_j}\), \(p'\in P^{\mathcal D_\ell}\), \(j\ne\ell\),
and \(f_s=f'_t\) for some \(s,t\).
\item[(CR10)] If \((C,D)\in R(r)\), \(r\sqsubseteq s\in\mathcal K\), and
\((C,D)\notin R(s)\), add \((C,D)\) to \(R(s)\).
\item[(CR11)] If \((C,D)\in R(r_1)\), \((D,E)\in R(r_2)\),
\(r_1\circ r_2\sqsubseteq r_3\in\mathcal K\), and
\((C,E)\notin R(r_3)\), add \((C,E)\) to \(R(r_3)\).
\end{enumerate}
The rules form a finite forward saturation.  CR1--CR7 and CR9--CR11
use the source Table~2 premises, with the typed CR11 update to \(R(r_3)\).
CR8 includes the feature-support condition just proved necessary.
The extra query nominal belongs to the same basic-description state as every
other nominal.  This remains the source's completion format; it introduces
neither a branching ABox tableau nor a blocking condition
\cite[table~2]{BaaderBrandtLutz2005}
\cite[table~2 and lem.~6]{BaaderBrandtLutz2005TR}.

\paragraph{Termination and the marked decision criterion.}
\label{sec:p09-elpp-query-correctness}
Put \(b=|\mathsf{BC}_{\mathcal K}|\) and \(m_R=|R_{\mathcal K}|\).
There are at most \(b(b+1)\) possible label memberships and \(m_Rb^2\)
possible role pairs.  Each successful rule application adds at least one
previously absent entry, so at most \(b(b+1)+m_Rb^2\) applications can succeed.
The rule instances can be searched in polynomial time: the sets are finite,
the reachability test is graph reachability, and the concrete satisfiability
and guarded implication tests are polynomial by \(p\)-admissibility.
Feature supports are read directly from finite labels.  Thus applying rules
until none is applicable terminates in polynomial time.

For the marked input \(\mathcal K\), let \(S,R\) be the resulting maps.
The decision criterion is
\[
 A_0\sqsubseteq_{\mathcal C}B_0
 \quad\Longleftrightarrow\quad
 B_0\in S(N_\ast)
 \ \lor\
 \bigl(\exists\{a\}\in\mathsf{BC}_{\mathcal K}\;
                         \bot\in S(\{a\})\bigr).
 \tag{D}
\]
The second disjunct includes \(N_\ast\).
It detects inconsistency of the marked CBox, which may arise either from
inconsistency of the original CBox or from impossibility of an
\(A_0\)-instance.  Both imply the queried subsumption.  Inconsistency forced
by a different, disconnected nominal must also be detected.

\paragraph{Two saturation runs.}
For a run using conjunction and role composition, take
\[
\begin{array}{rcl@{\qquad}rcl}
 A&\sqsubseteq&B, & A&\sqsubseteq&\exists r.C,\\
 C&\sqsubseteq&D, & \exists r.D&\sqsubseteq&E,\\
 B\sqcap E&\sqsubseteq&F, & r&\sqsubseteq&s,\\
 C&\sqsubseteq&\exists t.G, & s\circ t&\sqsubseteq&u.
\end{array}
\]
For the query \(A\sqsubseteq F\), adjoin \(N_\ast\sqsubseteq A\).
CR1 puts \(A,B\) in \(S(N_\ast)\) and \(D\) in \(S(C)\).
CR3 creates \((N_\ast,C)\in R(r)\) and \((C,G)\in R(t)\).
CR4 adds \(E\) to \(S(N_\ast)\), and CR2 adds \(F\).
CR10 also gives \((N_\ast,C)\in R(s)\), after which CR11 gives
\((N_\ast,G)\in R(u)\).  The first disjunct of (D) answers the query.

For the nominal case, take distinct concept names \(A,C,D,P,Q\), roles
\(r,s\), and the seven normalized GCIs
\[
\begin{gathered}
 A\sqsubseteq\exists r.C,\qquad A\sqsubseteq\exists s.D,\qquad
 C\sqsubseteq\{a\},\qquad D\sqsubseteq\{a\},\\
 C\sqsubseteq P,\qquad D\sqsubseteq Q,\qquad P\sqcap Q\sqsubseteq\bot.
\end{gathered}
\]
For \(A\sqsubseteq P\), marking gives \(A\in S(N_\ast)\) and the two
edges \((N_\ast,C)\in R(r)\), \((N_\ast,D)\in R(s)\).
Both successor labels contain \(\{a\}\); they contain \(P\) and \(Q\),
respectively.  Since \(D\) is reachable from the nominal \(N_\ast\),
\(C\rightsquigarrow D\) holds.  CR6 copies \(Q\) into \(S(C)\);
CR2 then adds \(\bot\) there and CR5 adds \(\bot\) to \(S(N_\ast)\).
Criterion (D) therefore answers yes.  Semantically, any \(A\)-instance would
have both successors equal to \(a^{\mathcal I}\), forcing that one element
into \(P^{\mathcal I}\cap Q^{\mathcal I}\).  Hence \(A\) is empty in every
model.  The original CBox nevertheless has a one-element model with all five
concept names and both roles empty.  This run checks the empty-query-concept
case as well as the nominal merge.  Both runs are handbook constructions.

\paragraph{Soundness of every completion rule.}
We prove the two invariants displayed above for every intermediate state,
relative to \(\mathcal K\), by induction on the sequence of additions.
Initialization is immediate: \(C\sqsubseteq_{\mathcal K}C,\top\), and
there are no role pairs.

For CR1, if \(x\in C^{\mathcal I}\), the induction hypothesis puts it in
\(C'^{\mathcal I}\), and the GCI puts it in \(D^{\mathcal I}\).
For CR2 the two hypotheses similarly put \(x\) in the intersection of
the two antecedents, so its GCI supplies the conclusion.
For CR3 the GCI supplies, for each such \(x\), an \(r\)-successor in
the indicated filler.  For CR4 the existing role-pair invariant supplies
\(y\in D^{\mathcal I}\) with \((x,y)\in r^{\mathcal I}\); its label
\(D'\) gives \(y\in D'^{\mathcal I}\), and the left-existential GCI
gives \(x\in E^{\mathcal I}\).
For CR5 the same witness would belong to \(\bot^{\mathcal I}\).
There can therefore be no \(x\in C^{\mathcal I}\), making the added
bottom subsumption sound.

For CR6, suppose \(x\in C^{\mathcal I}\).  The shared nominal label
gives \(x=a^{\mathcal I}\).
If the reachability sequence begins in \(C\), start with this \(x\);
if it begins in \(\{b\}\), start with \(b^{\mathcal I}\).
Apply the role-pair invariant successively along the finite sequence.
This produces an element \(y\in D^{\mathcal I}\), including the
zero-edge case.  The shared nominal label gives
\(y=a^{\mathcal I}=x\).  Every label copied from \(S(D)\) consequently
holds at \(x\).  The reachability premise remains necessary: membership
in a nominal alone does not ensure that an arbitrary concept has an instance.

For CR7, an element \(x\in C^{\mathcal I}\) would satisfy every positive
domain-\(j\) atom in \(S(C)\).  For each feature in \(F_j(C)\), its value
at \(x\) would thus be defined and lie in \(\Delta^{\mathcal D_j}\).
Those values form a solution of
\(\operatorname{con}_j(S(C))\), contradicting the rule's premise.
For CR8 the same assignment exists.  Its implication premise makes the
conclusion atom true under that assignment, and the support premise ensures
that every feature needed by that atom is among these defined values.
Thus the conclusion holds at \(x\) under the partial-feature semantics.
For CR9 a feature shared by positive atoms of two different domains would
have to take its single value in their disjoint carriers, which is impossible.

For CR10, the witness edge given by the role-pair invariant belongs to
the larger role by the RI.  For CR11, an instance \(x\) of \(C\) has an
\(r_1\)-successor \(y\) in \(D\); that particular instance \(y\) has an
\(r_2\)-successor \(z\) in \(E\).  The composition RI puts
\((x,z)\) in \(r_3^{\mathcal I}\).  This proves the second invariant.
All additions have now been covered.

If the first disjunct of (D) holds, soundness gives
\(N_\ast\sqsubseteq_{\mathcal K}B_0\), and (Q) gives the original query.
If its second disjunct holds, soundness would put the necessarily existing
element \(a^{\mathcal I}\) in \(\bot^{\mathcal I}\) in every
\(\mathcal K\)-model.  Thus \(\mathcal K\) is inconsistent and
\(N_\ast\sqsubseteq_{\mathcal K}B_0\) holds vacuously.  Equation (Q)
again yields the original query.

\paragraph{The reachable quotient for completeness.}
Assume that both disjuncts of (D) fail.
Let \(L\) be the set of basic descriptions reachable by zero or more
current role edges from some occurring nominal:
\[
 L=\{D\in\mathsf{BC}_{\mathcal K}:
       \exists\{a\}\in\mathsf{BC}_{\mathcal K}\;
       \{a\}\text{ has an }R\text{-path to }D\}.
\]
Here an \(R\)-path uses edges from any \(R(r)\).
Every nominal belongs to \(L\), in particular \(N_\ast\in L\).
If \(C\in L\) and \((C,D)\in R(r)\), then \(D\in L\) by appending
that edge.  Moreover \(\bot\notin S(C)\) for every \(C\in L\):
otherwise repeated CR5 applications backwards along a path to \(C\)
would put bottom in its initial nominal's label, contradicting the failed
second disjunct.

For any \(C\in\mathsf{BC}_{\mathcal K}\) and \(D\in L\), the relation
\(C\rightsquigarrow D\) holds, using a nominal-to-\(D\) path in its
definition.  Therefore, if \(C,D\in L\) share a nominal in their labels,
non-applicability of CR6 in both directions gives \(S(C)=S(D)\).
Define on \(L\)
\[
 C\sim D\quad\Longleftrightarrow\quad
 C=D\ \text{ or }\ 
 \exists\{a\}\in S(C)\cap S(D).
\]
Reflexivity and symmetry are immediate.  For transitivity, the only
nontrivial case has \(C,D\) sharing a nominal \(N\) and \(D,E\) sharing
another nominal.  The label equalities just proved give
\(S(C)=S(D)=S(E)\), so \(N\) also witnesses \(C\sim E\).
Thus \(\sim\) is an equivalence relation, and \(S(C)\) depends only on
the class \([C]\).

We also need agreement of outgoing role pairs:
\[
 C\sim C'\quad\Longrightarrow\quad
 \bigl((C,D)\in R(r)\ \Longleftrightarrow\
       (C',D)\in R(r)\bigr)
 \qquad(C,C'\in L).
 \tag{R}
\]
For one implication, induct on the first addition of the pair \((C,D)\)
during the finite run, keeping the final relation \(C\sim C'\) fixed.
If CR3 added the pair from a label \(E\) and
\(E\sqsubseteq\exists r.D\), final label agreement puts \(E\) in
\(S(C')\), so final saturation under CR3 gives \((C',D)\in R(r)\).
If CR10 added it, apply the induction hypothesis to the earlier pair in
the smaller role, then final CR10.  If CR11 added it using earlier pairs
\((C,E)\in R(r_1)\) and \((E,D)\in R(r_2)\), the induction hypothesis
transfers the first pair to \((C',E)\), and final CR11 supplies the result.
These are the only rules adding role pairs; the initial relations are empty.
Interchanging \(C,C'\) proves the converse.  This argument uses final
saturation, and does not require equivalence to have been present at the
earlier insertion stage.

\paragraph{Concrete witnesses on each quotient class.}
For \([C]\in L/{\sim}\) and domain \(j\), the feature set \(F_j(C)\)
and the conjunction \(\Gamma_j=\operatorname{con}_j(S(C))\) depend only
on the class.  Since \(\bot\notin S(C)\) and CR7 is no longer applicable,
\(\Gamma_j\) has a solution on \(F_j(C)\).
Let \(E_j\) be the finite set of basic domain-\(j\) atoms not in \(S(C)\)
all of whose features belong to \(F_j(C)\).
There is a solution of \(\Gamma_j\) making every member of \(E_j\)
false.  Indeed, if no such solution existed, every solution would satisfy
\(\bigvee E_j\).  If \(E_j\) is empty this contradicts satisfiability
directly; otherwise convexity says that \(\Gamma_j\) implies a member
of \(E_j\).  Its feature-support condition holds, so CR8 would add it,
contrary to final saturation.  Choose one such solution
\[
 \delta([C],j):F_j(C)\longrightarrow\Delta^{\mathcal D_j}.
 \tag{F}
\]
When \(F_j(C)\) is empty, the empty assignment suffices: all concrete
predicates have positive arity, so \(E_j\) is then empty.
There are only finitely many classes and domains, so this construction uses
only a finite choice of witnesses, and no choice of a representative of each
class.  By CR9 and absence of bottom, the sets \(F_j(C)\) for different
domains are pairwise disjoint.

Define an interpretation by
\[
\begin{aligned}
 \Delta^{\mathcal I}&=L/{\sim},\\
 A^{\mathcal I}&=\{[C]:A\in S(C)\}
       &&(A\in N_C\cap\mathsf{BC}_{\mathcal K}),\\
 a^{\mathcal I}&=[\{a\}]
       &&(\{a\}\in\mathsf{BC}_{\mathcal K}),\\
 r^{\mathcal I}&=\{([C],[D]):C,D\in L,\ (C,D)\in R(r)\}
       &&(r\in R_{\mathcal K}),\\
 f^{\mathcal I}([C])&=\delta([C],j)(f)
       &&(f\in F_j(C)).
\end{aligned}
\]
The domain is nonempty because it contains \([N_\ast]\).
Label agreement makes the concept and feature clauses independent of
representatives; disjoint feature supports make the indicated \(j\) unique.
Features are undefined elsewhere.  Concept and role names absent from the
finite vocabulary are interpreted as empty, and each individual name whose
nominal is absent is interpreted as \([N_\ast]\).  The role clause defines
a set of class pairs without choosing representatives; (R) permits the
source representative of any witnessing edge to be replaced as needed.

For every \(C\in L\) and \(D\in\mathsf{BC}_{\mathcal K}\cup\{\bot\}\), we have
\[
 [C]\in D^{\mathcal I}\quad\Longleftrightarrow\quad D\in S(C).
 \tag{T}
\]
For \(\top\) this follows from initialization; for \(\bot\) from its
empty extension and the absence of bottom in all retained labels.
For concept names it is the defining clause.
For a nominal \(D=\{a\}\), membership on the left means
\([C]=[\{a\}]\), which gives equal labels and hence
\(\{a\}\in S(C)\).  Conversely that label, together with
\(\{a\}\in S(\{a\})\), gives \(C\sim\{a\}\).

For a domain-\(j\) atom \(D=p(f_1,\ldots,f_k)\) in \(S(C)\),
every feature is in \(F_j(C)\), and (F) satisfies this conjunct, so the
semantic atom holds.  If the atom is absent and every feature lies in
\(F_j(C)\), it belongs to \(E_j\) and (F) makes it false.
If some feature does not lie in \(F_j(C)\), that feature is either
undefined at \([C]\) or takes a value in another domain.  The carriers
are disjoint, so either possibility makes the domain-\(j\) atom false.
This proves both directions of (T), including the partial-feature cases.

\paragraph{Verification of the normalized axioms and the query.}
For \(C_1\sqsubseteq D\), an instance \([E]\) of \(C_1\) has
\(C_1\in S(E)\) by (T).  CR1 then gives \(D\in S(E)\), and (T)
gives its conclusion.  For \(C_1\sqcap C_2\sqsubseteq D\), use both
labels and CR2 in the same way.
For \(C_1\sqsubseteq\exists r.C_2\), (T) and CR3 give
\((E,C_2)\in R(r)\).  The closure of \(L\) puts \(C_2\) in \(L\),
and its initial self-label makes \([C_2]\) an instance of \(C_2\).
This is the required successor.

For \(\exists r.C_1\sqsubseteq D\), take an \(r\)-edge from \([E]\)
to a class \([F]\) satisfying \(C_1\).
The role definition and (R) provide a representative \(F'\in[F]\)
with \((E,F')\in R(r)\).  By (T), \(C_1\in S(F')\);
CR4 gives \(D\in S(E)\), as required.
For \(r\sqsubseteq s\), every witnessing pair in \(R(r)\) belongs to
\(R(s)\) by CR10, so the semantic inclusion holds.
Finally suppose \([C]\,r_1^{\mathcal I}\,[D]\) and
\([D]\,r_2^{\mathcal I}\,[E]\).
Choose witnessing pairs \((C',D')\in R(r_1)\) and
\((D'',E')\in R(r_2)\), with the indicated equivalent representatives.
Since \(D'\sim D''\), (R) gives \((D',E')\in R(r_2)\).
CR11 then supplies \((C',E')\in R(r_3)\) for every
\(r_1\circ r_2\sqsubseteq r_3\) in \(\mathcal K\).
Thus every normalized GCI and RI holds in \(\mathcal I\).

The root class satisfies \(N_\ast\), while \(B_0\notin S(N_\ast)\)
makes it fail \(B_0\) by (T).  Hence \(\mathcal I\) is a
\(\mathcal K\)-countermodel to the right side of (Q), and its reduct
is a \(\mathcal C\)-countermodel to \(A_0\sqsubseteq B_0\).
This proves the remaining direction of (D).

\paragraph{Complexity and the other reasoning tasks.}
The concrete-atom prepass and NF schedule have the linear written symbol-size
bound just proved, and adjoining the query marker adds only a constant number
of name and axiom occurrences.  Their actual encoded output has polynomial
length in the encoded raw CBox and query, including concrete-predicate
parameters; the construction takes polynomial time in the original encoded
input length.
The finite vocabularies and every concrete constraint passed to a rule test
are bounded by this encoded output.  Thus \(p\)-admissibility and the
completion bound give polynomial cost in the original encoded input.
Consequently subsumption in the displayed
\(\mathcal{EL}^{++}(\mathcal D_1,\ldots,\mathcal D_n)\) language with
CBoxes is polynomial whenever the concrete domains are \(p\)-admissible.
All pairs of original concept names can be classified by polynomially many marked
runs.  This establishes the complexity conclusion corresponding to conference
Theorem~4 using the procedure and proof supplied here; the source locators for
its normalization, completion and model construction remain
\cite[lems.~1--3 and thm.~4]{BaaderBrandtLutz2005}
and \cite[\S3.1, lems.~5--7 and thm.~8]{BaaderBrandtLutz2005TR}.

For completeness of the task interface, an ABox is a finite set of
assertions \(C(a)\) and \(r(a,b)\).  Replace them, respectively, by the
GCIs \(\{a\}\sqsubseteq C\) and
\(\{a\}\sqsubseteq\exists r.\{b\}\).  Define
\(\mathcal C_{\mathcal A}\) to be the union of all these GCIs and the
original CBox \(\mathcal C\).
The singleton semantics makes each replacement equivalent to its assertion
in every interpretation; no unique-name assumption is added.
An ABox is consistent exactly when this combined CBox does not entail
\(\top\sqsubseteq\bot\), because abstract domains are nonempty.
Instance checking for \(C(a)\) is exactly
\(\{a\}\sqsubseteq_{\mathcal C_{\mathcal A}}C\).
Concept satisfiability is the failure of
\(C\sqsubseteq_{\mathcal C}\bot\).
These are polynomial-size reductions to the arbitrary-query normalization and
marked test already proved.  They give the four tasks listed in report
Theorem~8; that task list is broader than the literal conference Theorem~4.

\paragraph{Concrete examples and the convexity boundary.}
Proposition~5 gives two \(p\)-admissible source domains.  The rational domain
has unary \(\top_{\mathbb Q},=_{q},>_{q}\), binary equality, and binary
\(+_{q}\), where \((x,y)\in(+_q)^{\mathbb Q}\) iff \(x+q=y\).  The string
domain over the ISO~8859-1 alphabet has \(\top_{\mathbb S},=_{w}\), equality,
and \(\operatorname{conc}_{w}\), where \((x,y)\) belongs iff \(y=xw\).
The sources misprint the latter predicate's domain superscript as
\(\mathbb Q\); the surrounding string definition fixes the intended
\(\mathbb S\) reading.  By contrast, adding a rational \(\leq_q\) predicate,
or suitable prefix and suffix predicates to the string domain, gives the
paper's non-convex examples: a conjunction can force a disjunction without
forcing any individual disjunct
\cite[prop.~5]{BaaderBrandtLutz2005}
\cite[props.~10--11]{BaaderBrandtLutz2005TR}.

\paragraph{Exact failure boundaries.}
Section~4 returns to the smaller \(\mathcal{EL}\) grammar---concept names,
\(\top\), conjunction, and existential restriction---with general TBoxes.
The comparison-language names below have the following exact standing:
\(\mathcal{EL}^{(\neg)}\) permits complement only on concept names;
\(\mathcal{ELU}\) adds concept union; \(\mathcal{EL}^{\geq2}\) adds the
unqualified restriction ``at least two \(r\)-successors'';
\(\mathcal{EL}(\mathcal D)\) adds atoms from one concrete domain;
\(\mathcal{ELI}\) adds \(\exists r^{-}.C\);
\(\mathcal{EL}^{\leq1}\) adds ``at most one \(r\)-successor''; and
\(\mathcal{FL}_0\) has conjunction and universal value restriction.  With
these definitions, the literal conference results are:
\begin{center}
\small
\begin{tabularx}{\textwidth}{@{}l X X@{}}
\toprule
result & input and task & exact source strength\\
\midrule
Theorem 6 & \(\mathcal{EL}^{(\neg)}\), satisfiability and subsumption &
\(\mathsf{EXPTIME}\)-complete\\
Theorem 7 & \(\mathcal{ELU}\), subsumption &
\(\mathsf{EXPTIME}\)-complete\\
Theorem 8 & \(\mathcal{EL}^{\geq2}\), subsumption &
\(\mathsf{EXPTIME}\)-complete\\
Theorem 9 & \(\mathcal{EL}(\mathcal D)\) for every non-convex
\(\mathcal D\), subsumption & \(\mathsf{EXPTIME}\)-hard; not stated complete\\
Theorem 10 & \(\mathcal{ELI}\), subsumption & \(\mathsf{PSPACE}\)-hard;
the 2005 upper bound is \(\mathsf{EXPTIME}\), and the exact class is stated
open\\
Theorem 11 & \(\mathcal{EL}^{\leq1}\), subsumption &
\(\mathsf{EXPTIME}\)-complete\\
Theorem 12 & \(\mathcal{FL}_0\), subsumption &
\(\mathsf{EXPTIME}\)-complete\\
\bottomrule
\end{tabularx}
\end{center}

The reductions explain why these are constructor-sensitive results.  Theorem~6
replaces complex complements by fresh names constrained in both directions.
Theorem~7 replaces every atomic complement by a fresh name, forces each
original/fresh pair to cover \(\top\), and propagates a fresh failure concept.
Theorem~8 simulates a disjunctive choice with two existential witnesses: their
coincidence selects one branch and their distinctness, detected by
\((\geq2\ r)\), selects the other.  Theorem~9 uses a minimal disjunction forced
by non-convexity to simulate the single disjunctive GCI needed for hardness.
Theorem~10 translates \(A\sqsubseteq\forall r.B\) into
\(\exists r^{-}.A\sqsubseteq B\) and atomic complement into a propagated
failure concept.  Theorem~11 replaces left-hand value restrictions by
existentials and right-hand value restrictions by a conjunction of existence
and at-most-one, using total functional roles.  Theorem~12 unfolds an
\(\mathcal{FL}_0\) countermodel into the tree of finite role words, where every
role becomes a total function.  These are the source proof routes; cited
external hardness and upper-bound theorems are not re-proved by either paper
\cite[\S4, thms.~6--12]{BaaderBrandtLutz2005}
\cite[\S\S5--6]{BaaderBrandtLutz2005TR}.

The paragraph between Theorems~8 and~9 additionally reports
\(\mathsf{EXPTIME}\)-completeness when role terms admit, respectively,
relational complement in \(\Delta^{\mathcal I}\times\Delta^{\mathcal I}\),
relational union, or transitive closure; the report numbers this
three-language result Theorem~16.  Conversely, the paper's permitted RIs keep
a role chain only on the left and one role name on the right.  Its footnote and
report Theorem~12 warn that arbitrary role-value maps with a chain on the
right are undecidable.  Neither fact licenses transfer to another role
language.

Finally, \(\mathcal{EL}^{++}\) is the 2005 research logic.  OWL~2~EL is a
later W3C profile with its own functional-style grammar, datatype inventory,
axiom forms, and global restrictions.  The former influenced the latter; they
are not definitionally identical.  No ABox tableau, anonymous-node expansion,
blocking condition, OWL profile checker, or later role-regularity mechanism is
part of this saturation presentation.

\subsubsection{The DL-Lite query-rewriting presentation}
\label{sec:p09-dl-lite-query-rewriting}

\paragraph{Source and reconstruction boundary.}
The source is the complete forty-five-page author-hosted version of
the 2007 \emph{Journal of Automated Reasoning} article, printed pp.~385--429.
The language and model theory occur on pp.~389--395, the chase on
pp.~395--405, the universal-model argument on pp.~411--413, and
\textsc{PerfectRef} and its complexity analysis on pp.~413--425.  The source
prints a full journal argument, but several local omissions and index slips
must be repaired below rather than silently copied.  The later W3C
Recommendation is used only to delimit OWL~2~QL; it is not evidence for what
the 2007 authors called their logics \cite{CalvaneseEtAl2007}.

\paragraph{Three languages, with their categories kept distinct.}
Let \(A\) range over atomic concept names, \(P\) over atomic role names, and
\(a,b\) over individual constants.  The source grammar for
\(\mathrm{DL\mbox{-}Lite}_{core}\) is
\[
\begin{array}{rcl@{\qquad}rcl}
 B&::=&A\mid\exists R,&R&::=&P\mid P^{-},\\
 C&::=&B\mid\neg B,&E&::=&R\mid\neg R.
\end{array}
\]
Thus \(B,C,R,E\) stand respectively for basic concepts, general concepts,
basic roles, and general roles.  If \(R=P\), then \(R^{-}=P^{-}\), while
\((P^{-})^{-}=P\); double syntactic negations are removed.  A knowledge base is
\(\mathcal K=\langle\mathcal T,\mathcal A\rangle\), with finite TBox
\(\mathcal T\) and finite ABox \(\mathcal A\).  In the core language,
\(\mathcal T\) contains inclusions \(B\sqsubseteq C\), while \(\mathcal A\)
contains only \(A(a)\) and \(P(a,b)\).  In particular, the positive left side
of a concept inclusion is basic, and arbitrary general concepts are not
primitive ABox assertions.

\(\mathrm{DL\mbox{-}Lite}_{\mathcal R}\) adds inclusions \(R\sqsubseteq E\),
including positive and negative role inclusions.  It adds no functionality
axiom.  Separately, \(\mathrm{DL\mbox{-}Lite}_{\mathcal F}\) adds
\((\mathsf{funct}\ R)\) for a basic role \(R\), but no role inclusion.  The
paper studies these as two extensions of the core, not as one unnamed union.
It also proves that the former has the finite-model property while the latter
does not \cite[\S\S2.1--2.2]{CalvaneseEtAl2007}.

\paragraph{Interpretations, satisfaction, and the source's UNA.}
An interpretation \(I=(\Delta^I,\cdot^I)\) has nonempty domain and assigns
\[
\begin{aligned}
 A^I&\subseteq\Delta^I,&
 P^I&\subseteq\Delta^I\times\Delta^I,\\
 (P^{-})^I&=\{(y,x):(x,y)\in P^I\},&
 (\exists R)^I&=\{x:\exists y\,(x,y)\in R^I\},\\
 (\neg B)^I&=\Delta^I\setminus B^I,&
 (\neg R)^I&=(\Delta^I\times\Delta^I)\setminus R^I.
\end{aligned}
\]
Concept and role inclusions mean inclusion of the corresponding extensions;
\(A(a)\) and \(P(a,b)\) mean \(a^I\in A^I\) and
\((a^I,b^I)\in P^I\).  Functionality has the typed clause
\[
 (x,y),(x,z)\in R^I\Longrightarrow y=z.
\]
Because \(R\) may be inverse, this also expresses inverse functionality of an
atomic role.  The article expressly requires different constants to denote
different domain elements: its semantics includes the unique-name assumption
(UNA).  That hypothesis is used by the consistency and complexity development
and may not be transferred silently to OWL semantics
\cite[pp.~389--392]{CalvaneseEtAl2007}.

\paragraph{Queries, certain answers, and the two source complexity axes.}
A first-order query is an open formula \(q=\{\vec x\mid\varphi(\vec x)\}\),
with equality allowed.  A conjunctive query (CQ) has the form
\[
 \{\vec x\mid\exists\vec y\,
       \mathit{conj}(\vec x,\vec y)\},
\]
where the matrix is a conjunction of atoms and equalities; a union of
conjunctive queries (UCQ) is a finite disjunction of such matrices with one
free tuple.  The source states that equality atoms are eliminated by
identifying terms but does not spell out the distinct-constant failure case;
the following UNA-respecting completion is a handbook reconstruction.  For
each CQ, take the equivalence closure on its terms generated by the equality
atoms.  If one equivalence class contains two distinct constants, that
conjunction is false in every source interpretation and the disjunct is removed
from its UCQ.  Otherwise replace each class by its unique constant, if it has
one, and by one chosen representative variable if it does not; then delete the
equality atoms.  The empty union left when every disjunct is removed is the
identically false query.  Only the surviving equality-free CQs enter the
datalog notation and \textsc{PerfectRef}
\cite[pp.~392--393]{CalvaneseEtAl2007}.  Such a CQ over \(\mathcal K\) is
displayed as
\[
 q(\vec x')\leftarrow\mathit{conj}'(\vec x',\vec y'),
\]
using only atoms \(A(z)\) and \(P(z_1,z_2)\).  Head variables are
distinguished; body-only variables are non-distinguished.

Write \(\operatorname{Const}(\mathcal K)\) for the individual constants
occurring in \(\mathcal K\).  The certain answers are
\[
 \operatorname{ans}(q,\mathcal K)
 =\{\vec a\in\operatorname{Const}(\mathcal K)^{|\vec x|}:
    \vec a^{\,M}\in q^M
    \text{ for every }M\models\mathcal K\}.
\]
For an inconsistent \(\mathcal K\), the universal quantification has no
countermodel, so every tuple of the required arity is returned; the source
calls this finite set \(\operatorname{AllTup}(q,\mathcal K)\).

Only two complexity regimes are formally defined in this article.  In
\emph{data complexity}, \(\mathcal T\) and the query are fixed and only
\(|\mathcal A|\) varies.  In \emph{combined complexity}, every recognition
input varies---for query answering, the parameter is
\(|\mathcal K|+|Q|\).  The expressions \emph{taxonomic complexity} and
\emph{query complexity} belong to the later W3C table, not to this paper
\cite[\S\S2.3--2.4]{CalvaneseEtAl2007}.

\paragraph{The ABox database and first-order reducibility.}
For nonempty \(\mathcal A\), the finite database interpretation is
\[
\begin{aligned}
 \Delta^{db(\mathcal A)}&=\{a:a\text{ occurs in }\mathcal A\},&
 a^{db(\mathcal A)}&=a,\\
 A^{db(\mathcal A)}&=\{a:A(a)\in\mathcal A\},&
 P^{db(\mathcal A)}&=\{(a,b):P(a,b)\in\mathcal A\}.
\end{aligned}
\]
This is a minimal model of the ABox, not generally a model of the TBox.
Definition~3 says that UCQ answering is first-order reducible when, for every
fixed \(\mathcal T,Q\), one FO query \(Q_{\mathcal T}\), independent of
\(\mathcal A\), satisfies
\[
 \vec a\in\operatorname{ans}
   (Q,\langle\mathcal T,\mathcal A\rangle)
 \quad\Longleftrightarrow\quad
 \vec a\in Q_{\mathcal T}^{db(\mathcal A)}
\]
for every nonempty ABox and every active-domain tuple.  This is a uniform
definability statement.  A machine-space bound follows only after a fixed FO
formula is evaluated over the varying database; the two claims are not
synonyms.  Definition~3 says that the rewriting is over the alphabet of
\(\mathcal T\), while Section~5 assumes a query over \(\mathcal T\); the
printed phrase must not be read as licensing predicates in \(Q\) and then
forbidding them from its rewriting \cite[\S2.5]{CalvaneseEtAl2007}.

Every statement below whose right side contains \(db(\mathcal A)\) therefore
has the standing hypothesis \(\mathcal A\ne\varnothing\).  This is exactly the
input domain of the article's first-order-reducibility Definitions~2--3 and of
Theorems~41 and~46.  No empty-domain object \(db(\varnothing)\) is introduced.
When \(\mathcal A=\varnothing\), each TBox in either displayed variant has a
nonempty one-element model interpreting every atomic concept and role as empty;
query answering must be read directly from the query semantics, and a
non-Boolean answer has no active constants from which to form a tuple.

\paragraph{Positive inclusions and the chase.}
A \emph{positive inclusion} (PI) is \(B_1\sqsubseteq B_2\), or in the
\(\mathcal R\)-variant \(R_1\sqsubseteq R_2\); a \emph{negative inclusion}
(NI) has the same left category and a negated right side.  These names classify
existing TBox formulas and add no new object-language constructor.  For a
basic role \(R\), define the metalinguistic orientation abbreviation
\[
 \operatorname{ga}(R,a,b)=
 \begin{cases}
  P(a,b),&R=P,\\
  P(b,a),&R=P^{-}.
 \end{cases}
\]
It merely lets one chase rule cover a role and its inverse.

Put
\[
 \Gamma_{\mathcal A}:=
 \{a:a\text{ is an individual constant occurring in }\mathcal A\},
\]
and take a countably infinite fresh set \(\Gamma_N\) disjoint from
\(\Gamma_{\mathcal A}\).  Fix orders on the ABox facts, PIs, and
\(\Gamma_N\).  Starting with \(S_0=\mathcal A\), the chase adds one fact at a time by
the first applicable instance of
\[
\begin{array}{lll}
 \mathsf{cr1}:&A_1\sqsubseteq A_2,\ A_1(a)
     &\rightsquigarrow A_2(a),\\
 \mathsf{cr2}:&A\sqsubseteq\exists R,\ A(a)
     &\rightsquigarrow\operatorname{ga}(R,a,a_{new}),\\
 \mathsf{cr3}:&\exists R\sqsubseteq A,\ \operatorname{ga}(R,a,b)
     &\rightsquigarrow A(a),\\
 \mathsf{cr4}:&\exists R_1\sqsubseteq\exists R_2,\
       \operatorname{ga}(R_1,a,b)
     &\rightsquigarrow\operatorname{ga}(R_2,a,a_{new}),\\
 \mathsf{cr5}:&R_1\sqsubseteq R_2,\ \operatorname{ga}(R_1,a,b)
     &\rightsquigarrow\operatorname{ga}(R_2,a,b).
\end{array}
\]
Rules \(\mathsf{cr2}\) and \(\mathsf{cr4}\) fire only when the required outgoing
witness is absent; the other rules fire only when their conclusion is absent.
Rule \(\mathsf{cr5}\) belongs only to the \(\mathcal R\)-variant.  The printed
Definition~5 omits the saturated case.  The conservative completion used here
is
\[
 \text{if no PI applies to }S_j,\qquad S_{j+1}:=S_j.
\]
It makes every stage total without changing any generated fact or the intended
union \(\operatorname{chase}(\mathcal K)=\bigcup_{j\in\mathbb N}S_j\).
Proposition~6 proves fairness of the fixed-order construction
\cite[defs.~4--5 and prop.~6]{CalvaneseEtAl2007}.

Put \(\Gamma_C=\Gamma_{\mathcal A}\cup\Gamma_N\).  The canonical
interpretation reads the union chase literally:
\[
\begin{aligned}
 \Delta^{\operatorname{can}(\mathcal K)}&=\Gamma_C,&
 a^{\operatorname{can}(\mathcal K)}&=a\quad(a\in\Gamma_C),\\
 A^{\operatorname{can}(\mathcal K)}
   &=\{a:A(a)\in\operatorname{chase}(\mathcal K)\},&
 P^{\operatorname{can}(\mathcal K)}
   &=\{(a,b):P(a,b)\in\operatorname{chase}(\mathcal K)\}.
\end{aligned}
\]
Lemma~7 proves that it satisfies the ABox and every PI.  For a nonempty ABox,
Lemma~8 proves, in the \(\mathcal F\)-variant, that its functionality axioms
hold exactly when they already hold in \(db(\mathcal A)\).

\paragraph{Negative inclusions and finite consistency checking.}
The NI-closure \(\operatorname{cl}_n(\mathcal T)\) is a TBox obtained by the
following least closure.  It contains every NI in \(\mathcal T\), and, for the
\(\mathcal F\)-variant, every functionality assertion.  It is closed under:
\begin{enumerate}
\item if \(B_1\sqsubseteq B_2\in\mathcal T\) and either
\(B_2\sqsubseteq\neg B_3\) or \(B_3\sqsubseteq\neg B_2\) is already present,
add \(B_1\sqsubseteq\neg B_3\);
\item if \(R_1\sqsubseteq R_2\in\mathcal T\) and either
\(\exists R_2\sqsubseteq\neg B\) or \(B\sqsubseteq\neg\exists R_2\) is present,
add \(\exists R_1\sqsubseteq\neg B\);
\item under the same role premise, if either
\(\exists R_2^{-}\sqsubseteq\neg B\) or
\(B\sqsubseteq\neg\exists R_2^{-}\) is present, add
\(\exists R_1^{-}\sqsubseteq\neg B\);
\item under the same role premise, if either \(R_2\sqsubseteq\neg R_3\) or
\(R_3\sqsubseteq\neg R_2\) is present, add \(R_1\sqsubseteq\neg R_3\);
\item in the \(\mathcal F\)-variant, the presence of either
\(\exists R\sqsubseteq\neg\exists R\) or
\(\exists R^{-}\sqsubseteq\neg\exists R^{-}\) forces both; in the
\(\mathcal R\)-variant, the presence of either of those or
\(R\sqsubseteq\neg R\) forces all three.
\end{enumerate}
The \(\mathcal F\)-closure uses only the initial assertions, the first rule, and
its last \(\mathcal F\)-clause; the \(\mathcal R\)-closure uses the initial NIs,
the first rule, the three role-propagation rules, and the last
\(\mathcal R\)-clause.  These rules are metalinguistic closure rules, not
additions to the DL grammar.

For \(\mathcal A\ne\varnothing\), Lemmas~12 and~14 yield the decisive
equivalence
\[
 \mathcal K\text{ is satisfiable}
 \quad\Longleftrightarrow\quad
 \operatorname{can}(\mathcal K)\models\mathcal K
 \quad\Longleftrightarrow\quad
 db(\mathcal A)\models
   \langle\operatorname{cl}_n(\mathcal T),\mathcal A\rangle.
\]
For an assertion \(\alpha\) in the closure, the source turns its violation into
a Boolean FO formula \(\delta(\alpha)\):
\[
\begin{aligned}
 \delta((\mathsf{funct}\ P))
   &=\exists x\,y\,z\,[P(x,y)\land P(x,z)\land y\ne z],\\
 \delta((\mathsf{funct}\ P^{-}))
   &=\exists x\,y\,z\,[P(x,y)\land P(z,y)\land x\ne z],\\
 \delta(B_1\sqsubseteq\neg B_2)
   &=\exists x\,[\gamma_1(x)\land\gamma_2(x)],\\
 \delta(R_1\sqsubseteq\neg R_2)
   &=\exists x\,y\,[\rho_1(x,y)\land\rho_2(x,y)],
\end{aligned}
\]
where \(\gamma_i(x)\) is \(A_i(x)\), \(\exists y_iP_i(x,y_i)\), or
\(\exists y_iP_i(y_i,x)\) according as \(B_i\) is \(A_i\), \(\exists P_i\),
or \(\exists P_i^{-}\), and \(\rho_i(x,y)\) is \(P_i(x,y)\) or \(P_i(y,x)\)
according as \(R_i=P_i\) or \(R_i=P_i^{-}\).  Define
\[
 q_{\mathrm{unsat}}
 :=\bigvee_{\alpha\in\operatorname{cl}_n(\mathcal T)}\delta(\alpha),
\]
with the empty disjunction equal to \(\mathsf{false}\).  For
\(\mathcal A\ne\varnothing\), define the Boolean procedure used below by
\[
 \operatorname{Consistent}(\mathcal K)=\mathsf{true}
 \quad\Longleftrightarrow\quad
 q_{\mathrm{unsat}}^{db(\mathcal A)}=\mathsf{false}.
\]
Here \(q_{\mathrm{unsat}}\) is a metalinguistic Boolean first-order
inconsistency detector assembled from all closure violations, not the user's
UCQ \(Q\), and \(\operatorname{Consistent}\) is its truth-value test.  The two
abbreviations are needed later to state \(\operatorname{Answer}\) and to supply
the inconsistent-KB completion of Theorem~41; neither names a motive or an
object-language constructor.  Lemma~16 proves that evaluating this disjunction
gives a terminating and correct consistency test; Theorem~17 obtains
first-order reducibility of satisfiability
\cite[def.~9 and thms.~15--17]{CalvaneseEtAl2007}.

\paragraph{Homomorphism into every model and UCQ evaluation.}
For every satisfiable \(\mathcal R\)- or \(\mathcal F\)-knowledge base and every
model \(M\models\mathcal K\), Lemma~28 constructs
\[
 \psi:\Delta^{\operatorname{can}(\mathcal K)}\longrightarrow\Delta^M
\]
preserving atomic concept and role facts.  ABox symbols go to their denotations,
and each chase witness goes to a witness supplied by \(M\).  The printed proof
does not assign the unused members of \(\Gamma_N\); mapping all of them to one
fixed element of the nonempty \(\Delta^M\) makes \(\psi\) total without changing
any preserved fact.  The article prints Theorem~29 as
\(\operatorname{ans}(Q,\mathcal K)=Q^{\operatorname{can}(\mathcal K)}\).
The left side contains tuples of source constants, whereas an ordinary query
extension on the right may also contain fresh chase elements.  Its type-correct
content, for a UCQ \(Q\) of arity \(r\), is therefore
\[
 \operatorname{ans}(Q,\mathcal K)
 =\{\vec a\in\operatorname{Const}(\mathcal K)^r:
     \vec a^{\,\operatorname{can}(\mathcal K)}
       \in Q^{\operatorname{can}(\mathcal K)}\}.
\]
Theorem~31 distributes this set over the CQ disjuncts.  Its proof names
Theorem~28 at one point, but the equality it uses is Theorem~29; Lemma~28
alone supplies only the homomorphism \cite[\S4]{CalvaneseEtAl2007}.

\paragraph{Typed input to \textsc{PerfectRef}.}
In a query atom, an argument occurrence is \emph{bound} iff it is a
distinguished variable, a variable occurring at least twice in the body, or a
constant.  It is \emph{unbound} iff it is a body-only variable occurring once.
The printed underscore \(\_\) denotes such an occurrence; different underscores
are treated as different variables during unification.

A PI is applicable to \(A(x)\) exactly when its right side is \(A\).  It is
applicable to \(P(x_1,x_2)\) exactly when either \(x_2=\_\) and its right side is
\(\exists P\), or \(x_1=\_\) and its right side is \(\exists P^{-}\), or it is a
role inclusion whose right side is \(P\) or \(P^{-}\).  The first two side
conditions prevent a backward existential step from deleting a join, output
variable, or constant constraint.

For an applicable PI \(I\), Definition~32 defines the replacement
\(\operatorname{gr}(g,I)\) by all eleven cases:
\begin{center}
\scriptsize
\begin{tabularx}{\textwidth}{@{}X X X@{}}
\toprule
atom \(g\) & inclusion \(I\) & \(\operatorname{gr}(g,I)\)\\
\midrule
\(A(x)\) & \(A_1\sqsubseteq A\) & \(A_1(x)\)\\
\(A(x)\) & \(\exists P\sqsubseteq A\) & \(P(x,\_)\)\\
\(A(x)\) & \(\exists P^{-}\sqsubseteq A\) & \(P(\_,x)\)\\
\(P(x,\_)\) & \(A\sqsubseteq\exists P\) & \(A(x)\)\\
\(P(x,\_)\) & \(\exists P_1\sqsubseteq\exists P\) & \(P_1(x,\_)\)\\
\(P(x,\_)\) & \(\exists P_1^{-}\sqsubseteq\exists P\) & \(P_1(\_,x)\)\\
\(P(\_,x)\) & \(A\sqsubseteq\exists P^{-}\) & \(A(x)\)\\
\(P(\_,x)\) & \(\exists P_1\sqsubseteq\exists P^{-}\) & \(P_1(x,\_)\)\\
\(P(\_,x)\) & \(\exists P_1^{-}\sqsubseteq\exists P^{-}\) & \(P_1(\_,x)\)\\
\(P(x_1,x_2)\) & \(P_1\sqsubseteq P\) or \(P_1^{-}\sqsubseteq P^{-}\) &
\(P_1(x_1,x_2)\)\\
\(P(x_1,x_2)\) & \(P_1\sqsubseteq P^{-}\) or \(P_1^{-}\sqsubseteq P\) &
\(P_1(x_2,x_1)\)\\
\bottomrule
\end{tabularx}
\end{center}
The first nine rows rewrite concept inclusions; the last two are available only
with \(\mathcal R\)-role inclusions.

Write \(q[g/g']\) for replacing one selected atom.  If two atoms unify,
\(\operatorname{reduce}(q,g_1,g_2)\) applies a most general unifier to the whole
query and removes the duplicate atom.  Finally, \(\tau(q)\) turns each variable
that has become body-only and non-shared into an underscore.  This second pass
is required because a reduction can make a formerly shared variable unbound.

The exact Figure~2 control flow can now be stated without a new algorithmic
vocabulary.  Initialize \(PR:=\{q\}\).  At the start of each pass save
\(PR':=PR\).  For every \(q'\in PR'\), first add every applicable
\(q'[g/\operatorname{gr}(g,I)]\); then, for every unifiable atom pair, add
\(\tau(\operatorname{reduce}(q',g_1,g_2))\).  Stop when \(PR'=PR\) and return
\(PR\).  Newly added queries are processed on the next pass, and set union
suppresses duplicates \cite[def.~32 and fig.~2]{CalvaneseEtAl2007}.

\paragraph{Termination and its exact finite-state repair.}
Lemma~34 observes that no generated query has more body atoms than the input,
no new term other than an underscore is introduced, and the concept/role
vocabulary is finite.  If \(n\) is the source's query-size parameter and \(m\)
the number of concept and role names in \(\mathcal T\) or \(q\), the paper gives
\(m(n+1)^2\) possible atoms and prints
\((m(n+1)^2)^n\) as a coarse bound on queries.  Because the paper does not fix a
literal coding of heads and shorter bodies inside that power, a
convention-independent sufficient bound is obtained by putting
\(a=m(n+1)^2\):
\[
 (n+1)^n\sum_{i=0}^{n}a^i
 \leq (n+1)^{n+1}a^n\qquad(a\geq1).
\]
Quotienting by atom order, duplicate removal, or variable renaming only reduces
this finite number.  Each nonfinal pass adds a new set element, so termination
follows.  For fixed \(q\), the repaired bound remains polynomial in \(m\),
which is exactly what Lemma~42 later needs.

\paragraph{The source's student--tutor rewrite and answer.}
Example~35 uses
\[
\begin{aligned}
 \mathsf{Professor}&\sqsubseteq\exists\mathsf{TeachesTo},&
 \mathsf{Student}&\sqsubseteq\exists\mathsf{HasTutor},\\
 \exists\mathsf{TeachesTo}^{-}&\sqsubseteq\mathsf{Student},&
 \exists\mathsf{HasTutor}^{-}&\sqsubseteq\mathsf{Professor},\\
 \mathsf{Professor}&\sqsubseteq\neg\mathsf{Student}
\end{aligned}
\]
and begins with
\[
 q_0(x)\leftarrow
 \mathsf{TeachesTo}(x,y),\mathsf{HasTutor}(y,\_).
\]
Backward application and one reduction generate exactly
\[
\begin{array}{rcl}
 q_1(x)&\leftarrow&\mathsf{TeachesTo}(x,y),\mathsf{Student}(y),\\
 q_2(x)&\leftarrow&\mathsf{TeachesTo}(x,y),\mathsf{TeachesTo}(\_,y),\\
 q_3(x)&\leftarrow&\mathsf{TeachesTo}(x,\_),\\
 q_4(x)&\leftarrow&\mathsf{Professor}(x),\\
 q_5(x)&\leftarrow&\mathsf{HasTutor}(\_,x).
\end{array}
\]
The negative inclusion participates in consistency checking but never in
\textsc{PerfectRef}.  Example~38 evaluates this returned UCQ on
\[
 \{\mathsf{Student}(\mathsf{John}),
   \mathsf{HasTutor}(\mathsf{John},\mathsf{Mary}),
   \mathsf{TeachesTo}(\mathsf{Mary},\mathsf{Bill})\}.
\]
Queries \(q_2,q_3,q_5\) return Mary; the other three have no database match.
The union is therefore exactly \(\{\mathsf{Mary}\}\)
\cite[examples~35 and~38]{CalvaneseEtAl2007}.

\paragraph{Rewriting correctness.}
Lemma~39 states, for an \(\mathcal R\)-TBox, a CQ \(q\), its returned set \(PR\),
and every nonempty ABox for which \(\mathcal K\) is satisfiable,
\[
 \operatorname{ans}(q,\mathcal K)=PR^{db(\mathcal A)}.
\]
For soundness, each backward inclusion step maps any match of the rewritten
atom to a match of its predecessor, and reduction preserves a match of the
unreduced query.  Since every ABox fact occurs in the chase, iteration gives
that each active-domain tuple in \(PR^{db(\mathcal A)}\) belongs to the
constant-restricted canonical answer set displayed above, hence to
\(\operatorname{ans}(q,\mathcal K)\).

For completeness, take \(\vec a\in\operatorname{ans}(q,\mathcal K)\).  By the
type-correct form of Theorem~29, the corresponding canonical match is a finite
match of \(q\) in the union chase and already occurs at some finite stage.
Retain only the chase facts and parent applications needed for that match, and
reverse them.  Reversing a produced leaf is a Figure~2
backward-inclusion step; when two query atoms acquire one predecessor fact,
reduction unifies them.  At stage \(0\) the resulting query has an ABox match.
This gives the opposite inclusion.  The article expands the
\(\mathsf{cr2}\) case and declares the other four chase forms analogous.

Four printed slips are corrected here: the witness set is a subset of
\(\operatorname{chase}(\mathcal K)\), not of the interpretation object; the
reverse-step predecessor is \(\operatorname{chase}_{k-i}(\mathcal K)\), not
\(\operatorname{chase}_{i}(\mathcal K)\); the post-reduction cardinality is
\(|q_2|=|G_{k-i}|\), not \(|G_{k-i+1}|\); and Figure~2's literal operation is
\(\tau\circ\operatorname{reduce}\), even where the proof abbreviates it to
\(\operatorname{reduce}\) \cite[pp.~419--421]{CalvaneseEtAl2007}.

\paragraph{From one CQ to arbitrary consistent or inconsistent UCQs.}
For \(\mathcal A\ne\varnothing\), Figure~3's procedure has the typed value
\[
 \operatorname{Answer}(Q,\langle\mathcal T,\mathcal A\rangle)=
 \begin{cases}
  \operatorname{AllTup}(Q,\mathcal K),
       &\operatorname{Consistent}(\mathcal K)=\mathsf{false},\\
  \left(\displaystyle\bigcup_{q_i\in Q}
     \operatorname{PerfectRef}(q_i,\mathcal T)\right)^{db(\mathcal A)},
       &\text{otherwise}.
 \end{cases}
\]
Lemma~16 supplies the first branch, Lemmas~34 and~39 supply termination and
correctness of each CQ rewrite, and Theorem~31 supplies their union.  Theorem~40
therefore proves exact correctness for every \(\mathcal R\)-knowledge base in
this nonempty-ABox regime.

Theorem~41 says the result is first-order reducible, including inconsistent
knowledge bases.  Its satisfiable case is local.  For the inconsistent case,
the source describes a finite FO augmentation and refers to its reference~[10]
rather than printing the formula.  Using only the source definitions, the
missing formula is immediate: for answer arity \(r>0\), disjoin the returned
UCQ with
\[
 q_{\mathrm{unsat}}\land\bigwedge_{i=1}^{r}(x_i=x_i);
\]
for a Boolean query, use \(q_{\mathrm{unsat}}\) itself.  Because the domain of
\(db(\mathcal A)\) is exactly its active constants, this disjunct returns every
tuple exactly in the inconsistent case.  This is a handbook reconstruction of
the referred step, not a formula printed by the article.

\paragraph{Exact parameterized results and proof boundaries.}
The paper's conclusions are:
\begin{center}
\small
\begin{tabularx}{\textwidth}{@{}l X X@{}}
\toprule
result & hypotheses and varying input & exact conclusion and proof status\\
\midrule
Theorem 41 & fixed \(\mathcal R\)-TBox and fixed UCQ; arbitrary nonempty ABox &
UCQ answering is first-order reducible; inconsistent-case formula described
and referred, with the local completion above\\
Theorem 43 & fixed UCQ and fixed nonempty ABox for varying TBox; fixed UCQ and
TBox for varying nonempty ABox &
polynomial time in TBox size with the UCQ and ABox fixed, and a
\(\mathsf{LOGSPACE}\) data upper bound; not completeness claims\\
Theorem 44 & variable \(\mathcal R\)-KB, UCQ, and candidate tuple; size
\(|\mathcal K|+|Q|\) & \(\mathsf{NP}\)-complete combined complexity; local
nondeterministic rewrite sketch and hardness from CQ database evaluation with
empty TBox\\
Theorem 46 & fixed \(\mathcal F\)-TBox and fixed UCQ; arbitrary nonempty ABox &
first-order reducible, as an immediate transfer through removal of functionality
from the canonical chase\\
Theorem 47 & the three parameter regimes of Theorem~43 and Theorem~44, with a
fixed nonempty ABox in the TBox-size regime and a nonempty varying ABox in the
data regime &
polynomial in TBox size for fixed UCQ and ABox, \(\mathsf{LOGSPACE}\) in data, and
\(\mathsf{NP}\)-complete combined; stated by transfer, without a new proof body\\
Theorem 48 & fixed extended core TBox and fixed CQ; varying ABox &
\(\mathsf{coNP}\)-hard data complexity if one admits (i) \(\forall R.A\) or
\(\neg A\) on an inclusion left, (ii) either as a query atom, or (iii) concept
union on an inclusion right; hardness only, imported from reference~[19]\\
\bottomrule
\end{tabularx}
\end{center}
Lemma~42's polynomial TBox-size claim concerns the rewriting phase and keeps
the UCQ fixed; the answering-time TBox-size regimes above also fix the ABox.
Theorems~43 and~47 derive their data bound by evaluating a fixed rewriting and
fixed inconsistency query on \(db(\mathcal A)\).  No theorem says that
materializing all rewrites is
\(\mathsf{LOGSPACE}\) in the combined input; Theorem~44's
\(\mathsf{NP}\)-completeness rules out that conflation absent a class collapse
\cite[\S\S5.3--7]{CalvaneseEtAl2007}.

\paragraph{Why this source logic is not OWL~2~QL.}
The term \emph{OWL QL} does not occur in the 2007 article.  The later W3C
Recommendation says that OWL~2~QL is based on the DL-Lite family and that
\(\mathrm{DL\mbox{-}Lite}_{\mathcal R}\) supplies its logical underpinning; it
does not assert identity.  The nonidentity is formal: the article enforces UNA,
while OWL~2~QL does not; OWL~2~QL has its own structural grammar for individuals,
object and data properties, datatypes, axiom forms, and global restrictions;
and its normative complexity table includes axes and bounds not stated in this
paper.  Therefore \textsc{PerfectRef} is the 2007 research algorithm, not the
normative OWL~2~QL specification \cite[\S3]{W3COWL2Profiles2012}.

\subsubsection{W3C profiles as normative languages}

\paragraph{The dated normative bodies and their standing.}
The normative source is the W3C Recommendation
\emph{OWL 2 Web Ontology Language: Profiles (Second Edition)}, dated
11 December 2012.
The same-date Structural Specification and Direct Semantics bodies, needed to
type imported axioms, inherited restrictions, and inference tasks, are used
with it.
W3C makes these Recommendations publicly readable and citable under its
document-use terms, but they are not public-domain material or permission to
publish an altered substitute specification.  The handbook therefore preserves
the indispensable functional-syntax names and reconstructs the restrictions in
its own prose
\cite{W3COWL2Profiles2012,W3COWL2Syntax2012,W3COWL2DirectSemantics2012}.
Every CamelCase expression below is a W3C functional-syntax constructor or
axiom name.  The words class, datatype, object/data property expression, class
expression, data range, and axiom name the Recommendation's grammar categories;
none is a new handbook ontology object.

The Recommendation calls OWL~2~EL, OWL~2~QL, and OWL~2~RL \emph{profiles}:
independently defined structural fragments of OWL~2 chosen for different
reasoning regimes.  Profile membership is determined by the whole ontology's
functional-style grammar and global restrictions, even when the ontology is
serialized in another syntax.  None of the three profiles is a subset of
another; that does not mean that they are disjoint.

Let \(Ax\) denote the source's \emph{axiom closure}: the axioms in the import
closure, with anonymous individuals from distinct imported ontologies renamed
apart.  This is a metalinguistic input set, not an OWL entity or semantic
relation.  It is needed because declarations, entity typing, property
simplicity, role-chain regularity, and anonymous-individual conditions are
checked over \(Ax\), not over one displayed axiom in isolation
\cite[\S1]{W3COWL2Profiles2012}
\cite[\S\S3.4 and~11]{W3COWL2Syntax2012}.

\paragraph{Restrictions inherited by all three profiles.}
Before any profile-specific production is applied, every ontology must satisfy
the OWL~2~DL typing and structural conditions.  Classes and datatypes have
disjoint types; object, data, and annotation properties have pairwise disjoint
types; every class, datatype, and property occurrence must have the appropriate
declaration; individual declarations are optional.  A version IRI may occur
only when an ontology IRI occurs, and neither IRI may use reserved vocabulary
in OWL~2~DL.  Imported ontologies must themselves be OWL~2 ontologies, reserved
vocabulary may identify entities only in its specified built-in positions, and
a literal's lexical form must belong to its datatype's lexical space.

Every data range admitted by these profiles is unary.  Consequently, although
the main EL and RL productions inherit a one-or-more list of data properties
in \texttt{DataSomeValuesFrom} or \texttt{DataAllValuesFrom}, the inherited
arity equality leaves exactly one property in a well-formed profile ontology.
Likewise, the two property lists printed by \texttt{HasKey} may each be empty,
but at least one property must occur overall.

The remaining restrictions are stated over \(Ax\).
\texttt{owl:topDataProperty} may occur only as the superproperty in
\texttt{SubDataPropertyOf}.  Every datatype \(DT\) occurring in \(Ax\) must
satisfy exactly one of the following: \(DT\) is \texttt{rdfs:Literal}; \(DT\)
belongs to the OWL~2 datatype map; or \(DT\) is the subject of exactly one
\texttt{DatatypeDefinition} axiom in \(Ax\).  There must also be a strict
partial order on the occurring datatypes such that every datatype mentioned in
the defining data range lies below the datatype being defined.  This separate
order is the acyclicity condition.

A datatype introduced as the subject of \texttt{DatatypeDefinition} must also
be explicitly declared.  The Structural Specification gives every such
defined datatype an empty lexical space, so it cannot type a literal.  It also
supports no facets and therefore cannot occur as the datatype in
\texttt{DatatypeRestriction}.  The profile grammars already omit
\texttt{DatatypeRestriction}, but the literal prohibition remains substantive
because all three profiles admit literal-bearing forms
\cite[\S9.4]{W3COWL2Syntax2012}.

The property restrictions require their source definitions.  An
object-property expression \(OPE\) is either a named object property \(OP\) or
\texttt{ObjectInverseOf(\(OP\))}.  Define the metalinguistic involution
\[
 \operatorname{INV}(OP)=\texttt{ObjectInverseOf(\(OP\))},
 \qquad
 \operatorname{INV}(\texttt{ObjectInverseOf(\(OP\))})=OP.
\]
The set \(\operatorname{AllOPE}(Ax)\) is the smallest set containing \(OP\)
and \(\operatorname{INV}(OP)\) for every named object property \(OP\) occurring
in \(Ax\).  An \(OPE\) is \emph{composite in \(Ax\)} exactly when
\begin{enumerate}
\item it is \texttt{owl:topObjectProperty} or
\texttt{owl:bottomObjectProperty}; or
\item \(Ax\) contains a length-greater-than-one
\texttt{ObjectPropertyChain} inclusion whose target is \(OPE\) or
\(\operatorname{INV}(OPE)\); or
\item \(Ax\) contains \texttt{TransitiveObjectProperty(\(OPE\))} or
\texttt{TransitiveObjectProperty(\(\operatorname{INV}(OPE)\))}.
\end{enumerate}
Define \(\to\) to be the smallest relation on
\(\operatorname{AllOPE}(Ax)\) satisfying all five clauses:
\begin{enumerate}
\item \texttt{SubObjectPropertyOf(\(OPE_1\ OPE_2\))} gives
\(OPE_1\to OPE_2\);
\item \texttt{EquivalentObjectProperties(\(OPE_1\ OPE_2\))} gives both
directions;
\item \texttt{InverseObjectProperties(\(OPE_1\ OPE_2\))} gives
\(OPE_1\to\operatorname{INV}(OPE_2)\) and
\(\operatorname{INV}(OPE_2)\to OPE_1\);
\item \texttt{SymmetricObjectProperty(\(OPE\))} gives
\(OPE\to\operatorname{INV}(OPE)\); and
\item \(OPE_1\to OPE_2\) entails
\(\operatorname{INV}(OPE_1)\to\operatorname{INV}(OPE_2)\).
\end{enumerate}
Its reflexive-transitive closure is \(\to^{*}\).  The expression \(OPE\) is
\emph{simple in \(Ax\)} exactly when every \(OPE'\) with
\(OPE'\to^{*}OPE\) is noncomposite.  Cardinality restrictions and
\texttt{ObjectHasSelf}, together with functional, inverse-functional,
irreflexive, asymmetric, and disjoint object-property axioms, may mention only
simple object-property expressions.

Regularity requires the existence of a strict partial order \(<\) on
\(\operatorname{AllOPE}(Ax)\).  For the named object properties quantified by
the source it must satisfy
\[
 OP_1<OP_2
 \quad\Longleftrightarrow\quad
 \operatorname{INV}(OP_1)<OP_2,
\]
and, for all object-property expressions,
\[
 OPE_1<OPE_2\quad\Longrightarrow\quad
 \neg(OPE_2\to^{*}OPE_1).
\]
Each inclusion
\texttt{SubObjectPropertyOf(ObjectPropertyChain(\(OPE_1\ \cdots\ OPE_n\))
\(OPE\))}, with \(n\geq2\), must then have one of exactly five shapes:
\begin{enumerate}
\item \(OPE=\texttt{owl:topObjectProperty}\);
\item \(n=2\) and \(OPE_1=OPE_2=OPE\);
\item \(OPE_i<OPE\) for every \(1\leq i\leq n\);
\item \(OPE_1=OPE\) and \(OPE_i<OPE\) for every \(2\leq i\leq n\); or
\item \(OPE_n=OPE\) and \(OPE_i<OPE\) for every \(1\leq i<n\).
\end{enumerate}
The relations \(\to^{*}\) and \(<\) are not interchangeable: the former is
generated by the ontology and defines simplicity, whereas the latter is an
existential witness constrained not to reverse the former and governs legal
chains.  All of \(\operatorname{INV}\), \(\operatorname{AllOPE}\),
\emph{composite}, \(\to^{*}\), \emph{simple}, and \(<\) are source-defined
metalinguistic objects or predicates on \(Ax\), not OWL entities or new
object-language constructors; their later use in expression and axiom
restrictions is why each distinction is needed.

Finally, no anonymous individual may occur in same/different-individual or
negative-property assertions, or inside \texttt{ObjectOneOf} or
\texttt{ObjectHasValue}.  Define the \emph{anonymous-individual graph} of
\(Ax\) to have the anonymous individuals in \(Ax\) as vertices and an
undirected edge between \(\_\!:\!x\) and \(\_\!:\!y\) for each
\texttt{ObjectPropertyAssertion(\(OPE\ \_\!:\!x\ \_\!:\!y\))} in \(Ax\).
This graph must be a forest.  For each pair of anonymous individuals joined by
an edge of this graph, \(Ax\) may contain at most one object-property assertion
between the pair in either direction.  Every tree must contain at least one anonymous vertex having at
most one object-property assertion to a named individual, in either direction.
EL and QL exclude anonymous individuals, so these graph conditions are vacuous
there; RL admits them and inherits the conditions substantively
\cite[\S\S3.1, 5.8.1, and~11]{W3COWL2Syntax2012}.

The six admitted/excluded axiom shapes below are handbook-derived membership
checks against these productions and restrictions; none is a W3C-printed
example.

\paragraph{OWL~2~EL: one expression grammar and its axiom positions.}
EL admits named individuals, named object properties but no inverse
object-property expressions, and named data properties.  Its admitted
predefined datatype IRIs are
\begin{quote}\small
\texttt{rdf:PlainLiteral}, \texttt{rdf:XMLLiteral},
\texttt{rdfs:Literal}, \texttt{owl:real}, \texttt{owl:rational},
\texttt{xsd:decimal}, \texttt{xsd:integer},
\texttt{xsd:nonNegativeInteger}, \texttt{xsd:string},
\texttt{xsd:normalizedString}, \texttt{xsd:token}, \texttt{xsd:Name},
\texttt{xsd:NCName}, \texttt{xsd:NMTOKEN}, \texttt{xsd:hexBinary},
\texttt{xsd:base64Binary}, \texttt{xsd:anyURI}, \texttt{xsd:dateTime},
and \texttt{xsd:dateTimeStamp}.
\end{quote}
It expressly excludes the floating-point datatypes, the remaining signed and
fixed-width integer datatypes, \texttt{xsd:language}, and
\texttt{xsd:boolean}.  A nonreserved IRI may nevertheless name one declared,
acyclically defined custom datatype.

An EL data range is an admitted datatype, an intersection of at least two EL
data ranges, or a one-literal \texttt{DataOneOf}.  One class-expression grammar
serves every class position.  It contains a class, an intersection of at least
two EL expressions, a one-individual \texttt{ObjectOneOf},
\texttt{ObjectSomeValuesFrom}, \texttt{ObjectHasValue},
\texttt{ObjectHasSelf}, \texttt{DataSomeValuesFrom}, and
\texttt{DataHasValue}.  The object existential has an EL filler; the data
existential has the unary range just described.

The complete logical axiom inventory is:
\begin{center}
\small
\begin{tabularx}{\textwidth}{@{}l X@{}}
\toprule
family & admitted EL forms\\
\midrule
class axioms &
\texttt{SubClassOf}, \texttt{EquivalentClasses}, and
\texttt{DisjointClasses}, with EL expressions in every class position\\
object-property axioms &
ordinary or chained \texttt{SubObjectPropertyOf},
\texttt{EquivalentObjectProperties}, domain and range,
\texttt{ReflexiveObjectProperty}, and
\texttt{TransitiveObjectProperty}\\
data-property axioms &
subproperty, equivalence, domain, range, and
\texttt{FunctionalDataProperty}\\
datatype and key axioms &
\texttt{DatatypeDefinition} and \texttt{HasKey}; at least one object or data
key property must occur\\
assertions &
same/different individuals, class assertions, and positive or negative object
and data-property assertions, all with named individuals\\
\bottomrule
\end{tabularx}
\end{center}
Declarations and annotation axioms are also admitted.  Disjoint union,
disjoint object/data properties, functional object properties, inverses, and
the constructors absent from the displayed expression grammar are excluded
\cite[\S2.2 and app.~6.1]{W3COWL2Profiles2012}.

EL alone adds a range/chain condition.  The source says that \(Ax\)
\emph{imposes class-expression range \(CE\) on the named object property
\(OP_1\)} exactly when there are named object properties
\(OP_1,\ldots,OP_k\), for some \(k\geq1\), such that
\[
 \texttt{SubObjectPropertyOf(\(OP_i\ OP_{i+1}\))}\in Ax
 \quad(1\leq i<k)
\]
and
\[
 \texttt{ObjectPropertyRange(\(OP_k\ CE\))}\in Ax.
\]
This is a source-defined metalinguistic relation on \(Ax\), not an OWL axiom:
it differs from an explicit range axiom by following zero or more ordinary
subproperty inclusions.  That closure is needed in the following condition.
If \(Ax\) contains
\texttt{SubObjectPropertyOf(ObjectPropertyChain(\(OP_1\ \cdots\ OP_n\)) \(OP\))}
and imposes class-expression range \(CE\) on \(OP\), then \(Ax\) must also
impose \(CE\) on the final chain member \(OP_n\).  Hence
\begin{quote}\small\ttfamily
SubObjectPropertyOf(ObjectPropertyChain(:r :s) :t)\par
ObjectPropertyRange(:t :C)
\end{quote}
is excluded unless the closure also imposes range \texttt{:C} on
\texttt{:s}; adding \texttt{ObjectPropertyRange(:s :C)} is one direct repair.
The trigger need not be direct: replacing the second displayed axiom by
\texttt{SubObjectPropertyOf(:t :u)} together with
\texttt{ObjectPropertyRange(:u :C)} still makes \(Ax\) impose that range on
\texttt{:t}.

The axiom
\begin{quote}\small\ttfamily
SubClassOf(:A ObjectSomeValuesFrom(:r :B))
\end{quote}
is admitted by the EL grammar.  Replacing the right side by
\texttt{ObjectAllValuesFrom(:r :B)} is excluded because universal object
restriction is not an EL class expression.  These are grammar-membership
examples, not satisfiability claims
\cite[\S\S2.2.3 and~2.2.6]{W3COWL2Profiles2012}.

\paragraph{OWL~2~QL: different left and right class positions.}
QL admits named individuals, named data properties, and a named object property
or its inverse as an object-property expression.  Its predefined datatype
whitelist and blacklist are exactly EL's.  Its data ranges are admitted
datatypes and intersections of at least two QL data ranges; it excludes
\texttt{DataOneOf}, union, complement, and datatype restriction.

There is no freely substitutable QL class-expression grammar.  The two source
positions are:
\begin{center}
\small
\begin{tabularx}{\textwidth}{@{}l X@{}}
\toprule
position & admitted QL forms\\
\midrule
subclass expression &
a class; \texttt{ObjectSomeValuesFrom} with filler exactly
\texttt{owl:Thing}; or \texttt{DataSomeValuesFrom} with a QL data range\\
superclass expression &
a class; intersection of at least two superclass expressions; complement of a
subclass expression; \texttt{ObjectSomeValuesFrom} with an atomic class filler;
or \texttt{DataSomeValuesFrom} with a QL data range\\
\bottomrule
\end{tabularx}
\end{center}
The object property may be named or inverse in either existential position.
The main body inherits a one-or-more data-property list while Appendix~6.2
prints one; unary data ranges make exactly one the only well-formed case, so
the surface mismatch does not change profile membership.

The axiom grammar uses those positions exactly.  \texttt{SubClassOf} takes one
subclass expression on the left and one superclass expression on the right.
\texttt{EquivalentClasses} and \texttt{DisjointClasses} take two or more
subclass expressions; \texttt{DisjointUnion} is absent.  QL admits ordinary,
but not chained, object subproperties; object-property equivalence,
disjointness, and inverse axioms; object-property domain and range axioms whose
class position is a superclass expression; and reflexivity, symmetry, and
asymmetry.  Its data-property axioms comprise subproperty, equivalence,
disjointness, domain with a superclass expression, and range with a QL data
range.  \texttt{DatatypeDefinition} takes a QL data range.  Assertions comprise
\texttt{DifferentIndividuals}, \texttt{ClassAssertion} with an atomic
\texttt{Class}, and positive object- and data-property assertions.  Declarations
and annotation axioms are also admitted.  QL excludes functionality, inverse
functionality, transitivity, keys, \texttt{SameIndividual}, and negative-
property assertions.  The dated source is internally inconsistent about
irreflexivity: the overview calls it supported, while both operative grammar
locations omit it.  The reconstruction does not silently insert it
\cite[\S\S3.1--3.2 and app.~6.2]{W3COWL2Profiles2012}.

The following shape is admitted:
\begin{quote}\small\ttfamily
SubClassOf(ObjectSomeValuesFrom(:r owl:Thing)\par
\hspace*{2em}ObjectSomeValuesFrom(ObjectInverseOf(:s) :B))
\end{quote}
The following is excluded because a superclass existential requires an atomic
class filler:
\begin{quote}\small\ttfamily
SubClassOf(:A ObjectSomeValuesFrom(:r ObjectIntersectionOf(:B :C)))
\end{quote}
The Recommendation omits \texttt{SameIndividual} because its advertised
first-order-rewriting behavior is obtained without equality assertions; it
mentions equality-closure preprocessing in logarithmic space, but that does
not add the axiom to the profile grammar
\cite[\S\S3.1 and~3.2.3]{W3COWL2Profiles2012}.

\paragraph{OWL~2~RL: three class positions and inherited role restrictions.}
RL admits named and anonymous individuals subject to the inherited forest
conditions.  Object-property expressions may be named or inverse; data
properties are named.  It excludes the four predefined top and bottom object
and data properties.  Its admitted predefined datatype IRIs are
\begin{quote}\small
\texttt{rdf:PlainLiteral}, \texttt{rdf:XMLLiteral},
\texttt{rdfs:Literal}, \texttt{xsd:decimal}, \texttt{xsd:integer},
\texttt{xsd:nonNegativeInteger}, \texttt{xsd:nonPositiveInteger},
\texttt{xsd:positiveInteger}, \texttt{xsd:negativeInteger},
\texttt{xsd:long}, \texttt{xsd:int}, \texttt{xsd:short},
\texttt{xsd:byte}, \texttt{xsd:unsignedLong}, \texttt{xsd:unsignedInt},
\texttt{xsd:unsignedShort}, \texttt{xsd:unsignedByte},
\texttt{xsd:float}, \texttt{xsd:double}, \texttt{xsd:string},
\texttt{xsd:normalizedString}, \texttt{xsd:token},
\texttt{xsd:language}, \texttt{xsd:Name}, \texttt{xsd:NCName},
\texttt{xsd:NMTOKEN}, \texttt{xsd:boolean}, \texttt{xsd:hexBinary},
\texttt{xsd:base64Binary}, \texttt{xsd:anyURI}, \texttt{xsd:dateTime},
and \texttt{xsd:dateTimeStamp}.
\end{quote}
It expressly excludes \texttt{owl:real} and \texttt{owl:rational}; a declared
custom datatype may still have one acyclic legal definition.  An RL data range
is an admitted datatype or an intersection of at least two RL data ranges.

The three class positions are:
\begin{center}
\small
\begin{tabularx}{\textwidth}{@{}l X@{}}
\toprule
position & admitted RL forms\\
\midrule
subclass expression &
an atomic class other than \texttt{owl:Thing}; recursive intersection or union
of at least two subclass expressions; \texttt{ObjectOneOf} with one or more
individuals; \texttt{ObjectSomeValuesFrom};
\texttt{ObjectHasValue}; \texttt{DataSomeValuesFrom}; or
\texttt{DataHasValue}.  The object existential has a subclass-expression
filler, with a separately admitted \texttt{owl:Thing} filler\\
superclass expression &
an atomic class other than \texttt{owl:Thing}; recursive intersection of at
least two superclass expressions;
complement of a subclass expression; \texttt{ObjectAllValuesFrom};
\texttt{ObjectHasValue}; object maximum cardinality exactly \(0\) or \(1\),
either unqualified, qualified by a subclass expression, or explicitly qualified
by \texttt{owl:Thing}; \texttt{DataAllValuesFrom};
\texttt{DataHasValue}; or data maximum cardinality exactly \(0\) or \(1\),
optionally qualified by an RL data range\\
equivalence expression &
an atomic class other than \texttt{owl:Thing}; recursive intersection of at
least two equivalence expressions;
\texttt{ObjectHasValue}; or \texttt{DataHasValue}\\
\bottomrule
\end{tabularx}
\end{center}
The universal object filler is a superclass expression.  The object/data
maximum-cardinality productions have separately printed unqualified
alternatives; the object production additionally prints \texttt{owl:Thing}
separately because that class is not a subclass expression.  Anonymous
individuals are forbidden inside \texttt{ObjectOneOf} by the inherited global
restriction.  Every restricted object property must remain simple.

The axiom grammar maps those positions exactly.  \texttt{SubClassOf} takes a
subclass expression on the left and a superclass expression on the right.
\texttt{EquivalentClasses} takes two or more equivalence expressions;
\texttt{DisjointClasses} takes two or more subclass expressions; and
\texttt{DisjointUnion} is absent.  RL admits every Structural Specification
object-property axiom except reflexivity, including regular chains,
functionality, inverse functionality, irreflexivity, symmetry, asymmetry, and
transitivity; object-property domain and range take a superclass expression.
It admits every data-property axiom, with a superclass expression in domain and
an RL data range in range, and \texttt{DatatypeDefinition} likewise takes an RL
data range.  Every assertion form is admitted, with a superclass expression in
positive \texttt{ClassAssertion}.  \texttt{HasKey} takes a subclass expression
followed by its object- and data-property lists, with at least one property in
total.  Declarations and annotation axioms are also admitted.  RL's
functionality, cardinality, asymmetry, irreflexivity, and disjointness forms
remain subject to simplicity, and its chains remain subject to regularity
\cite[\S4.2 and app.~6.3]{W3COWL2Profiles2012}.

The following axiom is admitted:
\begin{quote}\small\ttfamily
SubClassOf(ObjectUnionOf(:A :B)\par
\hspace*{2em}ObjectAllValuesFrom(:r ObjectComplementOf(:C)))
\end{quote}
The following is excluded because an object existential is not an RL
superclass expression:
\begin{quote}\small\ttfamily
SubClassOf(:A ObjectSomeValuesFrom(:r :B))
\end{quote}

\paragraph{RL syntax is not the RL/RDF rule artifact.}
Section~4.2 defines the syntactic profile interpreted by Direct Semantics.
Section~4.3 separately gives a partial first-order axiomatization of
RDF-Based Semantics over a ternary triple predicate \(T\), organized into list,
equality, property, class, datatype, and schema-vocabulary tables.  It omits
RDF/RDFS axiomatic triples and some RDFS entailment rules.  For arbitrary RDF
graphs the Recommendation promises soundness of rule-derived entailments but
does not promise completeness.

Theorem~PR1 gives an equivalence only when both premise and conclusion are RL
ontologies, no IRI has more than one entity type, the premise has no annotation
subproperty/domain/range axioms, and every conclusion axiom is an atomic
named-individual class assertion, named object-property assertion, named
data-property assertion, or \texttt{SameIndividual}.  Only then is
Direct-Semantics entailment equivalent to first-order entailment from the
premise's RDF translation plus the RL/RDF rules.  The printed proof is a sketch
through depth-one normalization, DLP rules, and corresponding derivation trees.
It is not completeness for arbitrary RDF graphs, arbitrary conclusions,
annotations, or every RDF-Based-Semantics task
\cite[thm.~PR1]{W3COWL2Profiles2012}.

\paragraph{All six directed failures of profile inclusion.}
With the required declarations and absence of unrelated global violations
understood, the following one-axiom witnesses derive the source's non-nesting
claim:
\begin{center}
\small
\begin{tabularx}{\textwidth}{@{}l X X@{}}
\toprule
failed inclusion & witness in the first profile & exclusion from the second\\
\midrule
EL \(\nsubseteq\) QL & \texttt{TransitiveObjectProperty(:r)} &
QL omits transitivity\\
QL \(\nsubseteq\) EL & \texttt{InverseObjectProperties(:r :s)} &
EL omits inverse-property expressions and axioms\\
EL \(\nsubseteq\) RL & \texttt{ReflexiveObjectProperty(:r)} &
RL expressly omits reflexivity\\
RL \(\nsubseteq\) EL & \texttt{FunctionalObjectProperty(:r)}, with
\texttt{:r} simple & EL omits object-property functionality\\
QL \(\nsubseteq\) RL & \texttt{ReflexiveObjectProperty(:r)} &
QL's operative grammar admits reflexivity; RL excludes it\\
RL \(\nsubseteq\) QL & \texttt{TransitiveObjectProperty(:r)} &
RL admits transitivity; QL excludes it\\
\bottomrule
\end{tabularx}
\end{center}
The witnesses establish lack of inclusion, not disjointness
\cite[\S\S1--4]{W3COWL2Profiles2012}.

For the standardized OWL rows, both the reasoning task and the varying input
must be fixed before a complexity class is read.  Under Direct Semantics,
\emph{ontology consistency} asks whether the ontology has a model;
\emph{class-expression satisfiability} asks whether the checked expression has
a nonempty extension in some model of the ontology; \emph{class-expression
subsumption} asks whether the first checked expression's extension is included
in the second's in every model; and \emph{instance checking} asks whether a
named individual belongs to the checked expression's extension in every model.
A Boolean conjunctive query is a closed existential conjunction of class,
object-property, and data-property atoms, and query answering asks whether it
holds in every model.

The checked expressions have an additional Direct-Semantics input condition.
Every object-property expression used by a cardinality restriction or
\texttt{ObjectHasSelf} in the expression checked for satisfiability, in either
expression checked for subsumption, or in the expression checked for instance
membership must be simple in the ontology's axiom closure \(Ax\).  This is not
exhausted by the corresponding restriction on expressions already occurring
inside \(Ax\) \cite[\S2.5]{W3COWL2DirectSemantics2012}.

The four axes have the Recommendation's exact metalinguistic standing.  In
\emph{data complexity} only the total size of the assertions varies.  In
\emph{taxonomic complexity} only the total size of the nonassertional axioms
varies.  In \emph{query complexity} only the conjunctive query varies, so that
axis applies only to query answering.  In \emph{combined complexity} the
axioms and assertions, the conjunctive query for query answering, and the
expression or expressions being checked for class-expression satisfiability,
subsumption, or instance checking vary.  Components not named by the measure
are fixed.  These are input-size measures, not constructors of a description
logic \cite[\S5 and table~10]{W3COWL2Profiles2012}.

\paragraph{Reasoning-task table and proof ownership.}
The five Direct-Semantics reasoning tasks and four varying-input axes are
defined before Table~\ref{tab:p09-owl-profile-complexity}.  Those task meanings
govern the Direct-Semantics, profile, and OWL~1~DL rows.  The RDF-Based row is a
literal transcription of Recommendation Table~10; this presentation does not
reconstruct the separate RDF-Based-Semantics definitions of those tasks.  The
table reproduces every row, including RL, RDF-Based Semantics, OWL~1~DL,
membership-only bounds, lower bounds, open entries, conditional results, and
nonapplicability.  The Recommendation calls every entry worst-case.  It gives
no row-specific reduction, algorithm, proof, or technical-paper citation:
Section~5 defines the axes, sketches most class names, and records the table.
Consequently the table is authoritative for the dated classifications, but it
is not a local proof of them.  Theorem~PR1 is the Recommendation's separate,
narrow proof sketch; it does not prove Table~10.

The polynomial-size-NFA condition in Table~10 has no operational test,
construction, or size convention in Section~5.  The same section sketches
several complexity classes but omits definitions of \(\mathsf{EXPTIME}\) and
\(\mathsf{coNP}\), even though the table uses both.  These presentation gaps do
not refute a cell; they require the handbook to preserve its logical force
without manufacturing a missing definition or proof
\cite[\S5 and table~10]{W3COWL2Profiles2012}.

\paragraph{Printed grammar conflicts and the reconstruction policy.}
The dated Recommendation contains unresolved local defects:
\begin{enumerate}
\item The QL overview includes irreflexivity, but Section~3.2.5 and
Appendix~6.2 both omit it.
\item The QL overview both admits subclass-position
\texttt{DataSomeValuesFrom} and later lists it as unsupported.  The explicit
subclass production admits it and controls the reconstruction.
\item QL's main-body \texttt{DataSomeValuesFrom} production inherits one or
more data properties while Appendix~6.2 prints one.  Unary data ranges make one
the only well-formed profile case.
\item The RL object-property-axiom alternatives omit a separator before
\texttt{TransitiveObjectProperty}; both adjacent productions are independently
defined and the overview admits both.
\item RL Appendix~6.3 refers to undefined \texttt{superComplementOf} where the
main grammar defines \texttt{superObjectComplementOf}; it also leaves a
trailing separator after the sole \texttt{superDataMaxCardinality} alternative.
\item Section~4.2.3 calls \texttt{equivClassExpression} plural in one prose
sentence while every operative production uses the singular.
\end{enumerate}
The current W3C errata page contains no Profiles entry resolving these points.
The handbook therefore uses the unambiguous main productions, separately
defined alternatives, and inherited arity rules, and discloses every choice
rather than representing it as an official correction
\cite{W3COWL2Errata}.

\paragraph{Research lineages are influences, not profile identities.}
OWL~2~EL is based on the \(\mathcal{EL}^{++}\) line, but the profile has the OWL
entity/datatype system, functional-style axiom families, keys, assertions, and
inherited global restrictions rather than the 2005 paper's arbitrary
\(p\)-admissible concrete-domain parameter.  The 2008 update supplies the
final-chain range condition and shows why dropping it can make subsumption
undecidable; it is still not the normative profile grammar
\cite{BaaderBrandtLutz2005,BaaderBrandtLutz2008,W3COWL2Profiles2012}.

OWL~2~QL uses \(\mathrm{DL\mbox{-}Lite}_{\mathcal R}\) as its logical
underpinning, but it rejects the 2007 paper's required UNA, excludes
functionality and keys, and uses OWL position-sensitive syntax.  Conversely,
its grammar has OWL axiom names not present in the paper's three-variant
language.  The W3C \(\mathsf{AC}^{0}\) cells are later Recommendation claims,
not literal theorems of the 2007 article.

The 2006 SROIQ paper starts from the description-logical basis of OWL~1~DL,
adds the constructors and role restrictions leading toward OWL~2, and proves
decidability, but expressly leaves SROIQ's exact complexity open.  Direct
Semantics says that OWL~2 extends and corresponds to SROIQ where their
constructs overlap because the paper lacks OWL datatypes and punning.  Thus the
paper cannot own Table~10's later exact \(\mathsf{N2EXPTIME}\) classification
\cite{HorrocksKutzSattler2006,W3COWL2DirectSemantics2012}.

Finally, full structural OWL~2 under Direct Semantics and arbitrary RDF graphs
under RDF-Based Semantics are different semantic rows, not a fourth and fifth
profile.  Neither may be called \emph{OWL 2 Full} as though syntax, global
restrictions, and semantic regime were interchangeable.

\paragraph{The 2017 book presentation.}
The registered 2017 textbook separately contains its own EL completion,
DL-Lite rewriting, expressive-DL ladder, and OWL comparison in Chapters~6--8.
The chapter locators and scope are established, but those bodies are not part
of the available source basis.  Their exact algorithms, examples, theorem
statements, and proofs are not reconstructed from the two research papers or
the W3C Recommendation \cite{BaaderHorrocksLutzSattler2017}.  Those three
presentations fulfill only their own source roles.
